% concatenated from main.tex
\documentclass{article}

% if you need to pass options to natbib, use, e.g.:
    \PassOptionsToPackage{numbers, sort&compress}{natbib}
% before loading neurips_2026

% The authors should use one of these tracks.
% Before accepting by the NeurIPS conference, select one of the options below.
% 0. "default" for submission
%\usepackage{neurips_2026}
% the "default" option is equal to the "main" option, which is used for the Main Track with double-blind reviewing.
% 1. "main" option is used for the Main Track
%  \usepackage[main]{neurips_2026}
% 2. "position" option is used for the Position Paper Track
%  \usepackage[position]{neurips_2026}
% 3. "eandd" option is used for the Evaluations & Datasets Track
 % \usepackage[eandd]{neurips_2026}
 % if you want to opt-in for a de-anomymized submission (not recommended, but OK):
 %\usepackage[eandd, nonanonymous]{neurips_2026}
% 4. "creativeai" option is used for the Creative AI Track
%  \usepackage[creativeai]{neurips_2026}
% 5. "sglblindworkshop" option is used for the Workshop with single-blind reviewing
 % \usepackage[sglblindworkshop]{neurips_2026}
% 6. "dblblindworkshop" option is used for the Workshop with double-blind reviewing
%  \usepackage[dblblindworkshop]{neurips_2026}`

% After being accepted, the authors should add "final" behind the track to compile a camera-ready version.
% 1. Main Track
 % \usepackage[main, final]{neurips_2026}
% 2. Position Paper Track
%  \usepackage[position, final]{neurips_2026}
% 3. Evaluations & Datasets Track
 % \usepackage[eandd, final]{neurips_2026}
% 4. Creative AI Track
%  \usepackage[creativeai, final]{neurips_2026}
% 5. Workshop with single-blind reviewing
%  \usepackage[sglblindworkshop, final]{neurips_2026}
% 6. Workshop with double-blind reviewing
%  \usepackage[dblblindworkshop, final]{neurips_2026}
% Note. For the workshop paper template, both \title{} and \workshoptitle{} are required, with the former indicating the paper title shown in the title and the latter indicating the workshop title displayed in the footnote.
% For workshops (5., 6.), the authors should add the name of the workshop, "\workshoptitle" command is used to set the workshop title.
% \workshoptitle{WORKSHOP TITLE}

% "preprint" option is used for arXiv or other preprint submissions
\usepackage[preprint]{neurips_2026}

% to avoid loading the natbib package, add option nonatbib:
%    \usepackage[nonatbib]{neurips_2026}

\usepackage[utf8]{inputenc} % allow utf-8 input
\usepackage[T1]{fontenc}    % use 8-bit T1 fonts
\usepackage{hyperref}       % hyperlinks
\usepackage{url}            % simple URL typesetting
\usepackage{booktabs}       % professional-quality tables
\usepackage{amsfonts}       % blackboard math symbols
\usepackage{nicefrac}       % compact symbols for 1/2, etc.
\usepackage{microtype}      % microtypography
\usepackage{xcolor}         % colors
\usepackage{adjustbox}
% Additional packages for this paper.
\usepackage{amsmath}
\usepackage{amssymb}
\usepackage{mathtools}
\usepackage{amsthm}
\usepackage{bm}
\usepackage{graphicx}
\usepackage{float}
\usepackage{subcaption}
\usepackage{array}
\usepackage{multirow}
\usepackage{tcolorbox}
\usepackage{algorithm}
\usepackage{algorithmic}

\usepackage[capitalize,noabbrev]{cleveref}
\tcbuselibrary{skins,breakable}

\graphicspath{{fig/}}

\newcommand{\R}{\mathbb{R}}
\newcommand{\E}{\mathbb{E}}
\newcommand{\Ent}{\mathrm{H}}
\newcommand{\I}{\mathrm{I}}
\newcommand{\KL}{\mathrm{KL}}
\newcommand{\Tr}{\mathrm{Tr}}
\newcommand{\dd}{\mathrm{d}}
\newcommand{\vx}{\bm{x}}
\newcommand{\vz}{\bm{z}}
\newcommand{\vv}{\bm{v}}
\newcommand{\vu}{\bm{u}}
\newcommand{\veps}{\bm{\epsilon}}
\newcommand{\marg}{\mathrm{marg}}
\newcommand{\cond}{\mathrm{cond}}
\newcommand{\cm}{\mathrm{cm}}
\newcommand{\bcr}{\mathrm{BCR}}
\newcommand{\diver}{\nabla\!\cdot}
\newcommand{\score}{\nabla \log}
\newcolumntype{P}[1]{>{\raggedright\arraybackslash}p{#1}}

\newcommand{\inlineeqtag}[1]{\refstepcounter{equation}\label{#1}\nobreak\hfill\textup{(\theequation)}\par}

\providecommand{\edmfidcell}[2]{%
  \mbox{#1\,{\tiny $\pm\,#2$}}%
}
\providecommand{\edmbestfidcell}[2]{%
  \begingroup\setlength{\fboxsep}{0.8pt}\colorbox{yellow!30}{\edmfidcell{#1}{#2}}\endgroup%
}
% \providecommand{\alphaplddtcell}[2]{%
%   \mbox{#1\,{\tiny $\pm\,#2$}}%
% }
  \providecommand{\alphaplddtcell}[2]{%
    \mbox{#1\,{\fontsize{3.5pt}
  {5pt}\selectfont $\pm\,#2$}}%
  }
% \providecommand{\alphaplddtcellbold}[2]{%
%   \mbox{\textbf{#1}\,{\tiny $\pm\,#2$}}%
% }
  \providecommand{\alphaplddtcellbold}[2]{%
    \mbox{\textbf{#1}\,{\fontsize{3.5pt}
  {5pt}\selectfont $\pm\,#2$}}%
  }
\providecommand{\alphabestplddtcell}[2]{%
  \begingroup\setlength{\fboxsep}{0.8pt}\colorbox{yellow!30}{\alphaplddtcell{#1}{#2}}\endgroup%
}
\providecommand{\alphabestplddtcellbold}[2]{%
  \begingroup\setlength{\fboxsep}{0.8pt}\colorbox{yellow!30}{\alphaplddtcellbold{#1}{#2}}\endgroup%
}


\newcolumntype{P}[1]{>{\raggedright\arraybackslash}p{#1}}

\theoremstyle{plain}
\newtheorem{theorem}{Theorem}
\newtheorem{proposition}{Proposition}
\newtheorem{lemma}{Lemma}
\theoremstyle{definition}
\newtheorem{assumption}{Assumption}
\theoremstyle{remark}
\newtheorem{remark}{Remark}

\newtcolorbox{computebox}{
  enhanced,
  breakable,
  colback=white,
  colframe=white,
  boxrule=0pt,
  borderline={0.5pt}{0pt}{gray!55,dotted},
  arc=1.5pt,
  left=5pt,
  right=5pt,
  top=4pt,
  bottom=4pt
}

\tcolorboxenvironment{proposition}{
  enhanced,
  breakable,
  colback=gray!7,
  colframe=gray!45,
  boxrule=0.5pt,
  arc=1.5pt,
  left=5pt,
  right=5pt,
  top=4pt,
  bottom=4pt
}

\tcolorboxenvironment{assumption}{
  enhanced,
  breakable,
  colback=white,
  colframe=black,
  boxrule=0.35pt,
  arc=2pt,
  borderline west={2pt}{0pt}{gray!45},
  left=6pt,
  right=5pt,
  top=4pt,
  bottom=4pt
}

% Note. For the workshop paper template, both \title{} and \workshoptitle{} are required, with the former indicating the paper title shown in the title and the latter indicating the workshop title displayed in the footnote.

% \title{Entropy Across the Bridge: Conditional--Marginal Entropy Rates for Flow and Schr\"odinger Samplers}

\title{Entropy Across the Bridge: Conditional--Marginal Discretization for Flow and Schr\"odinger Samplers}

\author{%
  \normalfont\mdseries
  \textbf{Bruno Trentini}$^{1,2}$\thanks{Correspondence to: \texttt{brunod@nvidia.com}, \texttt{luca.ambrogioni@donders.ru.nl}} \qquad
  \textbf{Dejan Stancevic}$^{3}$ \qquad
  \textbf{Michael M.~Bronstein}$^{2,4}$ \\[0.5em]
  \textbf{Alexander Tong}$^{4}$ \qquad
  \textbf{Luca Ambrogioni}$^{3*}$ \\[1em]
  \small
  $^{1}$NVIDIA Corporation, Santa Clara, CA, USA \\
  $^{2}$University of Oxford, Dept.\ of Computer Science, Oxford, UK \\
  $^{3}$Donders Institute for Brain, Cognition, and Behaviour, Radboud University, Nijmegen, NL \\
  $^{4}$AITHYRA, Research Institute for Biomedical AI, Vienna, AT%
}

% \title{Entropy Across the Bridge: Conditional--Marginal Entropy Rates for Low-Atep Generative Sampling}
% \author{%
%   \normalfont\mdseries
%   \begin{tabular}[t]{ccc}
%     \textbf{Bruno Trentini}$^{1,2}$\thanks{Correspondence to: \texttt{brunod@nvidia.com}, \texttt{luca.ambrogioni@donders.ru.nl}} &
%     \textbf{Dejan Stancevic}$^{3}$ &
%     \textbf{Michael M.~Bronstein}$^{2,4}$ \\[1.5em]
%     \shortstack{\textbf{Alexander Tong}$^{4}$} &
%     {} &
%     \shortstack{\textbf{Luca Ambrogioni}$^{3}$} \\
%   \end{tabular}\\[1em]%
%   \small%
%   $^1$ NVIDIA Corporation, Santa Clara, CA, USA \quad%
%   $^2$ University of Oxford, Oxford, UK \\%
%   $^3$ Radboud University, Nijmegen, NL \quad%
%   $^4$ AITHYRA, Vienna, AT%
% }
% The \author macro works with any number of authors. There are two commands
% used to separate the names and addresses of multiple authors: \And and \AND.
%
% Using \And between authors leaves it to LaTeX to determine where to break the
% lines. Using \AND forces a line break at that point. So, if LaTeX puts 3 of 4
% authors names on the first line, and the last on the second line, try using
% \AND instead of \And before the third author name.


% \author{%
%   Bruno Trentini$^{1,2}$ \quad
%   Dejan Stancevic$^{3}$ \quad
%   Michael M. Bronstein$^{2,4}$ \\
%   Alexander Tong$^{4}$ \quad
%   Luca Ambrogioni$^{3}$ \\
%   \\
%   $^{1}$NVIDIA \\
%   $^{2}$University of Oxford, Department of Computer Science, Oxford, UK \\
%   $^{3}$Donders Institute for Brain, Cognition and Behaviour, Radboud University, Nijmegen, The Netherlands \\
%   $^{4}$AITHYRA, Research Institute for Biomedical Artificial Intelligence, Vienna, Austria \\
%   % examples of more authors
%   % \And
%   % Coauthor \\
%   % Affiliation \\
%   % Address \\
%   % \texttt{email} \\
%   % \AND
%   % Coauthor \\
%   % Affiliation \\
%   % Address \\
%   % \texttt{email} \\
%   % \And
%   % Coauthor \\
%   % Affiliation \\
%   % Address \\
%   % \texttt{email} \\
%   % \And
%   % Coauthor \\
%   % Affiliation \\
%   % Address \\
%   % \texttt{email} \\
% }
% \author{%
%   Bruno Trentini \\
%   NVIDIA, University of Oxford \\
%   Oxford, UK \\
%   \And
%   Dejan Stancevic \\
%   Donders Institute for Brain, Cognition and Behaviour, Radboud University \\
%   Nijmegen, The Netherlands \\
%   \And
%   Michael M.~Bronstein \\
%   University of Oxford, AITHYRA \\
%   Oxford, UK; Vienna, Austria \\
%   \AND
%   Alexander Tong \\
%   AITHYRA \\
%   Vienna, Austria \\
%   \And
%   Luca Ambrogioni \\
%   Donders Institute for Brain, Cognition and Behaviour, Radboud University \\
%   Nijmegen, The Netherlands \\
% }

% \author{%
%   Bruno Trentini$^{1,2}$ \quad
%   Dejan Stancevic$^{3}$ \quad
%   Michael M.~Bronstein$^{2,4}$ \\
%   Alexander Tong$^{4}$ \quad
%   Luca Ambrogioni$^{3}$ \\
%   \\[0.75em]
%   $^{1}$NVIDIA \\
%   $^{2}$University of Oxford, Department of Computer Science, Oxford, UK \\
%   $^{3}$Donders Institute for Brain, Cognition and Behaviour, Radboud University, Nijmegen, The Netherlands \\
%   $^{4}$AITHYRA, Research Institute for Biomedical Artificial Intelligence, Vienna, Austria \\
% }

\begin{document}


\maketitle


% \begin{abstract}
% Flow and Schr\"odinger bridge samplers are often deployed under small inference budgets, where only a few neural network function evaluations (NFE) are available and the time grid can dominate sample quality. Current practice mostly uses fixed grids or heuristics such as linear, cosine, sigmoid, power, or log-SNR schedules, while learned scheduling methods optimize the grid empirically. Recent entropy-based schedules give a principled signal for diffusion models, but they do not separate the conditional and marginal geometry that defines a bridge. We propose a conditional--marginal entropy-rate objective for bridge-aware discretization. For Gaussian Brownian bridges, the conditional term has a closed form and predicts a symmetric U-shaped profile; in general, we estimate the matched conditional and marginal divergence terms with Hutchinson traces, allocate steps by inverse CDF, and temper the rate when the raw profile is too concentrated. EDM and AlphaFlow serve different roles: EDM is a diffusion-style image sampler that tests robustness under model mismatch, while AlphaFlow is a protein flow closer to the bridge setting. On EDM, log-tempered entropy gives the best tested five-step ODE FID, $184.9\pm4.8$, which is $6.8\%$ lower than linear and $21.1\%$ lower than cosine. On AlphaFlow, the conditional--marginal estimate recovers the U-shaped profile predicted by the bridge calculation; entropic log1p achieves endpoint pLDDT $86.3$, $88.3$, and $89.3$ at 5, 10, and 25 steps, compared with $86.1$, $88.1$, and $89.0$ for linear. We show that conditional--marginal entropy rate is a bridge-aware inference-time signal for low-NFE sampling, provided it is estimated as the two-term bridge quantity and calibrated into a stable numerical grid.
% \end{abstract}

\begin{abstract}
    For a fixed flow-based generative model under a small inference budget, sample quality can depend strongly on where the sampler spends its few function evaluations. Flow matching and Schr\"odinger bridges define probability paths, yet their inference grids are usually heuristic or inherited from one-endpoint diffusion. We derive a conditional--marginal entropy-rate objective for bridge-aware discretization, separating endpoint-conditioned bridge geometry from marginal flow evolution, and use it to build a training-free entropic inference-time scheduler from first principles. For Gaussian Brownian bridges this rate is closed-form and U-shaped, motivating boundary-heavy nonuniform grids. On trained two-dimensional bridge/flow models, the estimated profile recovers the predicted shape and improves 10-step ODE-Heun MMD over linear by 18.1\%, with a paired 22.7\% SDE-Heun improvement in the same low-NFE sweep. On EDM/CIFAR-10, the entropic time-discretization gives the best tested five-step FID ($186.3\pm4.0$ versus $200.5\pm2.9$ for linear and $238.0\pm5.3$ for cosine). On AlphaFlow protein generation, entropic conditional--marginal (cond-marg) scheduling shows advantage in low-NFE regimes on both CAMEO22 and ATLAS benchmarks. These results support entropy-rate scheduling as a practical low-budget allocation signal for high-dimensional bridge and flow samplers.


    % This is especially relevant in high-dimensional generation, where extra steps are costly and heuristic grids remain standard. We introduce a principled bridge-aware scheduler based on a conditional--marginal entropy rate signal, separating endpoint-conditioned bridge geometry from marginal flow evolution. For Gaussian Brownian bridges, the conditional component is boundary-heavy with a mid-time minimum, motivating non-uniform step placement. On EDM/CIFAR-10, the log-tempered entropy proxy yields the best tested five-step ODE-Heun Inception FID ($186.3\pm4.0$ versus $200.5\pm2.9$ for linear and $238.0\pm5.3$ for cosine). alphaf generation, raw Entropic conditioinal--marginal scheduling is competitive with the strongest five-NFE endpoint-confidence baselines. These results support entropy-rate scheduling as a practical low-budget alloication signal for high-dimensional bridge and flow samplers.

    % We introduce a bridge-aware scheduler based on a conditional--marginal entropy-rate signal. The construction separates the endpoint-conditioned geometry of a bridge from the marginal evolution of the flow. For Gaussian Brownian bridges, the conditional component has boundary singularities and a mid-time minimum, motivating boundary-aware step placement. On trained two-dimensional bridge/flow models, estimated entropy profiles recover this shape, and the resulting scheduler improves low-step MMD by up to 18.1\% at 10 steps over linear. We then test the principle outside the controlled setting. On EDM/CIFAR-10, where bridge assumptions are mismatched, raw entropy over-concetrates the grid. Log-tempered schedules gives the best tested five-step ODE-Heun Inception FID $186.3\pm4.0$ versus $200.5\pm2.9$ for linear and $238.0\pm5.3$ for cosine. On AlphaFlow protein generation, where structure generation and folding-style evaluation are expensive, raw entropic scheduling is competitive with the strongest five-NFE endpoint-confidence baselines, while larger-NFE plDDT exposes finite-grid limits. These results support conditional--marginal entropy as a practical low-budget allocation signal for bridge and flow samplers, not as a direct endpoint-quality guarantee

    % where the conditional term captures endpoint-conditioned bridge geometry, and subtracting the marginal term removes volume changed shared by the learned flow.

    % Low-budget generative sampling is often limited not by the model alone, but by where a sampler spends its limited function evaluations. Diffusion models, flow matching, and Schro\"odinger bridges all trace probability paths, yet their inference grids are usually fixed by heuristics or by entropy schedules that treat the path as a one-endpoint process. We introduce a bridge-aware scheduling mechanism based on a conditional--marginal entropy-rate signal. . In Gaussian Brownian Bridges, this quantity has a closed-form profile with boundary singularities and a mid-time minimum. We test our {scheduler? estimator?} in two high-dimensional settings with different roles: on EDM/CIFAR-10,a diffusion-model mismatch case, raw entropy over-concentrates the grid, while log-tempered scheduling gives the best tested five-step ODE-Heun Inception FID, $186.3\pm4.0$, compared with $200.5\pm2.9$ for linear and $238.0\pm5.3$ for cosine. On AlphaFlow protein generation, a bridge-like flow setting, the conditional--marginal diagnostic recovers boundary-heavy allocation showing that raw Entropic scheduling is competitive with the strongest five-NFE endpoint-confidence baseliines, while larger-NFR plDDT exposes finite-ggrid and metric-specific limits. Together, these results support conditional--marginal entropy as a practical allocation signal for low-budget bridge and flow sampling, not as a universal guarantee of endpoint quality.
\end{abstract}
% \begin{abstract}
% Flow and Schr\"odinger bridge samplers are often deployed under small inference budgets, where only a few neural network function evaluations (NFE) are available and the time grid can dominate sample quality. Current practice mostly uses fixed grids or heuristics such as linear, cosine, sigmoid, power, or log-SNR schedules, while learned scheduling methods optimize the grid empirically. Recent entropy-based schedules give a principled signal for diffusion models, but they do not separate the conditional and marginal geometry that defines a bridge. We propose a conditional--marginal entropy-rate objective for bridge-aware discretization. For Gaussian Brownian bridges, the conditional term has a closed form and predicts a symmetric U-shaped profile; in general, we estimate the matched conditional and marginal divergence terms with Hutchinson traces, allocate steps by inverse CDF, and temper the rate when the raw profile is too concentrated. We test our method on EDM and AlphaFlow, serving different roles: EDM is a diffusion-style image sampler that tests robustness under model mismatch, while AlphaFlow is a protein flow closer to the bridge setting. On EDM, log-tempered entropy gives the best tested five-step ODE FID, $184.9\pm4.8$, which is $6.8\%$ lower than linear and $21.1\%$ lower than cosine. {\color{red} To be updated: On AlphaFlow, the conditional–marginal estimate closely recovers the U-shaped profile predicted by the bridge calculation, achieving endpoint pLDDT consistently higher than the linear schedule across low-NFE settings, with the advantage becoming slightly more pronounced at larger step counts. These results support the claimt that the conditional–marginal entropy rate can serve as a bridge-aware inference-time signal for low-NFE sampling, provided it is estimated using the full two-term bridge quantity and calibrated onto a numerically stable grid.}
% \end{abstract}


\section{Introduction}
\label{sec:introduction}

Modern generative samplers often define a continuous-time path and then approximate it with a small number of discrete evaluations. Diffusion models, flow matching, stochastic interpolants, and Schr\"odinger bridges all share this numerical bottleneck \citep{ho2020ddpm,kingma2021vdm,karras2022edm,lipman2023flow,albergo2023building,debortoli2021dsb,liu2023i2sb,tong2024simulation}. We use NFE\footnote{NFE counts neural network calls during sampling. It is a useful proxy for latency, energy use, and serving cost because these calls dominate inference in many diffusion, flow, and bridge samplers.} to denote the number of neural network function evaluations performed during sampling. At high NFE, many grids become similar because integration error is small. In the low-NFE regime, a misplaced step can dominate the behavior of the sampler. This issue is no longer only about image synthesis. The same sampling budget appears in protein design, molecular generation, materials discovery, biological trajectory modeling, and other AI-for-science systems where each extra model call can be expensive.

% I want here to make the reader feel the pain of low step sampling. Like the model can be very good but if the few time points are badly placed the output can be bad anyway. This is very informational but maybe useful because it says why this is a global AI problem and not only a small scheduler trick. Need later make it less clunky.
% Maybe phrase: we are not changing the model, we are asking where the model should spend its few chances to look at the path.

The standard response is to choose a useful time grid. Linear, cosine, sigmoid, power, and log-SNR grids encode different beliefs about where denoising or transport is difficult \citep{nichol2021improved,kingma2021vdm,karras2022edm,lin2024common}. Higher-order solvers and predictor-corrector schemes change the update rule \citep{lu2022dpm,lu2022dpmsolverpp,zhao2023unipc,zhang2022deis,tan2026stork}. Optimized or learned schedulers search for better timesteps or trajectory parameterizations, including Align Your Steps, optimized-step rules, dynamic-transport schedules, and recent few-step scheduler learning \citep{sabour2024alignsteps,xue2024optimized,tsimpos2025lipschitz,min2026bezierflow}. Align Your Steps optimizes a discrete schedule for diffusion sampling, Lipschitz or transport-based rules use smoothness proxies, and learned scheduler families fit parameterized curves. These methods show that the grid matters, albeit they do not by themselves say which bridge quantity should be measured. Scheduling and distillation also solve different problems: distillation changes the model weights, while the present method changes the inference grid and can therefore be combined with a distilled sampler. The distinction matters because the bridge and flow literature has moved beyond one-endpoint image diffusion. Schr\"odinger bridges now appear in simulation-free flow and score matching, image translation, energy-based sampling, biological trajectories, discrete spaces, and structured domains \citep{tong2024simulation,liu2023i2sb,liu2025asbs,tamogashev2026dataenergy,ksenofontov2025categorical,tang2026branched,wyrwal2026topological}. Flow matching has also become central in protein and molecular generation, including AlphaFlow, Proteina, BioEmu, MolFlow, FlowMM, and Wasserstein flow matching \citep{jing2024alphafoldmeetsflowmatching,proteina2025,bioemu2024,irwin2024molflow,miller2024flowmm,haviv2025wasserstein}. In these settings, a sampler is not just reversing a noising process because it transports between structured endpoints, often under geometric constraints.

Recent entropic time work gives a principled schedule for diffusion models by using conditional entropy as a time coordinate \citep{stancevic2025entropic,stancevic2026information}. A Schr\"odinger bridge carries an additional layer: conditional paths indexed by endpoint information and a marginal population flow obtained after those conditions are mixed. A single-field entropy story is therefore too small for the bridge, because the rate must compare how volume changes along the paired conditional path with how volume changes after endpoint mixing. The bridge story needs a two-term object. Let $Z$ denote the bridge condition, such as a data endpoint or endpoint pair, and let $X_t$ denote the state at time $t$. Under a smooth probability-flow description, we show that
\begin{equation}
  \frac{\dd}{\dd t}\Ent(Z\mid X_t)
  =
  \E_{Z,X_t\mid Z}\!\left[\diver \vv_t(X_t\mid Z)\right]
  -
  \E_{X_t}\!\left[\diver \bar{\vv}_t(X_t)\right].
  \label{eq:intro_identity}
\end{equation}
The first term measures conditional volume change along paired bridge paths. The second measures marginal volume change after the bridge conditions have been mixed. Their difference is the conditional--marginal (cond--marg) information-rate signal studied here as a bridge-aware scheduling signal. The proposed grid is obtained by normalizing the magnitude of this signal and inverting its cumulative distribution.
% I should maybe explain this equation more slowly.....
% Extra equation idea, not for final:
% \begin{equation}
%   \text{bridge information speed}
%   \approx
%   \text{conditional volume change}
%   -
%   \text{mixed population volume change}.
% \end{equation}
% \begin{figure}[t]
%   \centering
%   \includegraphics[width=0.88\linewidth]{entropic_condmarg_variants_comparison.pdf}
%   \caption{Conditional--marginal scheduler variants used to connect the Brownian-bridge calculation to the empirical grids. The raw rate preserves the bridge signal but can over-concentrate near endpoints; the log1p transform is the regularized default used in the high-dimensional experiments.}
%   \label{fig:theory_overview}
% \end{figure}
    % \begin{figure}[t]
    %   \centering
    %   \includegraphics[width=0.92\linewidth]{entropy_tau_grid.pdf}
    %   \caption{Entropy-rate profiles and induced inverse-CDF time grids in the controlled transport setting. The figure shows the core mechanism used throughout the paper: measure where the path carries information, normalize that rate, and place time nodes by the corresponding cumulative distribution.}
    %   \label{fig:entropy_tau_grid_hero}
    % \end{figure}

% \begin{figure}[t]
%   \centering
%   \begin{minipage}[t]{0.49\linewidth}
%     \centering
%     \adjustbox{max width=\linewidth, max height=0.32\textheight, valign=t}{%
%       \includegraphics{entropy_tau_grid.pdf}%
%     }
%   \end{minipage}
%   \hfill
%   \begin{minipage}[t]{0.49\linewidth}
%     \centering
%     \adjustbox{max width=\linewidth, max height=0.32\textheight, valign=t}{%
%       \includegraphics{fig_transport_scenarios_grid_tarchic_vector_field_rose_points.pdf}%
%     }
%   \end{minipage}
%   \caption{Controlled transport setup and entropy-induced grids. Left: entropy-rate profiles and induced inverse-CDF time grids show the core mechanism: measure where the path carries information, normalize that rate, and place time nodes by the corresponding cumulative distribution. Right: synthetic transport scenarios used as controlled probes, contrasting continuous--continuous, continuous--discrete, discrete--continuous, and discrete--discrete endpoint structure.}
%   \label{fig:entropy_tau_grid_hero}
% \end{figure}
    \begin{figure}[t]
      \centering
      \includegraphics[width=\linewidth]{edm_alphaflow_entropy_profiles_tarchic.pdf}
      \caption{Entropy profiles for EDM and AlphaFlow. EDM is sharply endpoint-concentrated and requires log tempering; AlphaFlow recovers the boundary-heavy cond--marg profile predicted by the Brownian-bridge calculation. Shaded bands mark BCR endpoint windows.}
      \label{fig:entropy_curves}
    \end{figure}\textbf{}
This paper makes five contributions: (a) an entropy-rate criterion for inference-time discretization in flow and bridge samplers; (b) a cond--marg estimator that implements \Cref{eq:intro_identity}; (c) a Brownian-bridge closed form showing why the conditional term is U-shaped and singular near the endpoints; (d) a link between entropy-weighted grids and local ODE error; (e) a staged empirical path from controlled two-dimensional bridge/flow models to high-dimensional samplers.

%The two-dimensional models are the pre-flight check that asks whether learned vector fields recover the theoretical U-shape rate, whether the induced grids help in the low-NFE regime, and whether ODE and SDE samplers tell the same qualitative story before the estimator is moved to expensive domains. EDM and AlphaFlow then play different roles. EDM is not a bridge model and is not trained on the entropy objective; therefore, it is a useful \textit{mismatch case} for asking whether a tempered information-rate grid can still help a pretrained diffusion-based image sampler \citep{karras2022edm}. AlphaFlow is the closer high-dimensional test for the bridge diagnostic because its flow matching sampler exposes the conditional and marginal fields needed to estimate the cond--marg rate. There we ask whether the cond--marg rate itself shows the U-shaped structure predicted by the Brownian bridge calculation \citep{jing2024alphafoldmeetsflowmatching}. Results and Discussion then separate schedule geometry from endpoint metrics, because the high-dimensional models tested in this work were pretrained and were not optimized with the entropy criterion.

% EDM and AlphaFlow are intentionally different stress tests. EDM is not a Schr\"odinger bridge, so it tests whether entropy can still help when the raw profile must be tempered. AlphaFlow is closer to a bridge because it is a protein flow-matching model, so it tests whether the conditional--marginal rate has the U-shaped structure predicted by the Brownian bridge calculation. The endpoint metrics are interpreted separately from allocation diagnostics because AlphaFlow was trained with its original objective, not the entropy objective studied here.

%LUCA: Removed to save space
%\begin{table}[t]
%  \caption{Relation to the closest entropy-based scheduling work. The key change is the bridge-specific contrast between conditional and marginal divergence.}
%  \label{tab:stancevic}
%  \centering
%  \small
%  \setlength{\tabcolsep}{3pt}
%  \begin{tabular}{@{}P{0.14\linewidth}@{\hspace{0.02\linewidth}}P{0.39\linewidth}@{\hspace{0.02\linewidth}}P{0.39\linewidth}@{}}
%    \toprule
%    Aspect & Entropic diffusion time & Conditional--marginal bridge time \\
%    \midrule
%    Object & One-endpoint diffusion process & Conditional bridge or flow path \\
%    Signal & Entropy of data given noised state & $\dot{\Ent}(Z\mid X_t)$ from conditional minus marginal divergence \\
%    Estimator & Diffusion loss or single-field formula & Matched conditional and marginal divergence estimates \\
%    Shape & Often asymmetric in noise time & Symmetric U-shape for Brownian bridges \\
%    Use & Inference-time reparameterization & Inference-time grid for bridge and flow samplers \\
%    \bottomrule
%  \end{tabular}
%\end{table}

\section{Preliminaries}
\label{sec:setup}

\textbf{Probability flows.} Let $p_0$ be a data distribution on $\R^d$ and $p_1$ be a reference distribution. A deterministic probability flow is a density path $(p_t)_{t\in[0,1]}$ and vector field $\bar{\vv}_t$ satisfying
\begin{equation}
  \partial_t p_t(\vx) + \diver\big(p_t(\vx)\bar{\vv}_t(\vx)\big)=0.
  \label{eq:continuity}
\end{equation}
Sampling integrates the learned field from the reference endpoint toward data on a decreasing time grid $1=t_0>t_1>\cdots>t_N=0$. This numerical grid is the object we study.

\textbf{Flow matching.} Flow matching often starts from a conditional interpolation. Given a condition $Z$, define a conditional law $p_t(\vx\mid Z)$ and a conditional vector field $\vv_t(\vx\mid Z)$ satisfying
\begin{equation}
  \partial_t p_t(\vx\mid Z) + \diver\big(p_t(\vx\mid Z)\vv_t(\vx\mid Z)\big)=0.
  \label{eq:conditional_continuity}
\end{equation}
The marginal field is the posterior average
%\begin{equation}
  $\bar{\vv}_t(\vx)
  =
  \E\!\left[\vv_t(\vx\mid Z)\mid X_t=\vx\right]$,
%  \label{eq:marginal_field}
%\end{equation}
whenever the conditional and marginal descriptions are compatible. This is the population field targeted by standard flow matching losses \citep{lipman2023flow,liu2023rectified,albergo2023building}.

\textbf{Optimal transport and Schr\"odinger bridges.} Optimal transport asks for a low-cost coupling between two endpoint distributions. A Schr\"odinger bridge adds stochasticity by selecting the path measure closest in relative entropy to a reference process while matching the same endpoints \citep{schrodinger1931,schrodinger1932,leonard2014schrodinger}. In the Brownian reference case, the bridge path measure can be written as a mixture of Brownian bridges weighted by an entropic optimal transport plan:
\begin{equation}
  P((X_t)_{t\in[0,1]})
  =
  \int Q((X_t)_{t\in[0,1]}\mid x_0,x_1)
  \,\dd \Pi_{2\sigma^2}(x_0,x_1).
  \label{eq:sb_mixture}
\end{equation}
For this case the natural condition is $Z=(X_0,X_1)$. The distinction between the conditional field and the marginal field is not cosmetic. It is the mathematical expression of the bridge constraint.

\textbf{Time grids and NFE.} Given a nonnegative rate $r(t)$, an entropic grid is constructed through
\begin{equation}
  Q(t)=\frac{\int_0^t r(s)\,\dd s}{\int_0^1 r(s)\,\dd s},
  \qquad
  t_k=Q^{-1}\!\left(\frac{k}{N}\right).
  \label{eq:inverse_cdf}
\end{equation}
The number of intervals, or equivalently the number of model evaluations for a one-step solver, is the NFE budget. The central question is what rate should be used for a bridge. The rate derived below is $r(t)=|\frac{\dd}{\dd t}\Ent(Z\mid X_t)|$, optionally transformed by $\log(1+r)$ when the raw signal is too concentrated near singular endpoints. Expanded derivations for the entropy identities, Brownian bridge field, and ODE allocation rule are given in \Cref{app:full_derivations,app:bb_details,app:sde_ode_details,app:ode_error}.

The empirical sections are guided by three questions summarized in \Cref{app:hypotheses}. First, does the Brownian-bridge calculation predict a U-shaped conditional rate, and do small learned two-dimensional bridge probes show the same boundary-heavy pattern? Second, does the cond--marginal estimator reveal bridge-specific allocation in a high-dimensional model? Third, when the rate is converted into a grid, does it help low-NFE sampling without being confused with endpoint quality or model training effects?

% Second, do learned two-dimensional bridge/flow models preserve that shape once the vector field is trained rather than analytic? Third, do ODE and SDE sampler respond consistently to the resulting grids, and why should the high-dimensional analysis focus on the deterministic probability-flow ODE? Fourth, does the conditional--marginal estimator reveal bridge-specific allocation in high-dimensional models without confusing grid geometry with endpoint quality or training objective effects?

% Second, does the conditional--marginal estimator reveal bridge-specific allocation in a high-dimensional model? Third, when the rate is converted into a grid, does it help low-NFE sampling without being confused with endpoint quality or model training effects?

\section{Conditional--marginal entropy rate}
\label{sec:theory}

The bridge constraint asks for more than the entropy of the current marginal state. A conditional bridge path can contract because it is resolving a particular endpoint pair, while the population flow can expand or contract differently after those endpoint pairs are averaged. A schedule based on only one of these effects confounds path geometry with population geometry. The derivation below isolates the difference. It uses the standard regularity needed to differentiate entropy through a continuity equation.

\begin{assumption}
\label{ass:smooth}
All conditional and marginal densities are smooth and positive on their support, the relevant vector fields are continuously differentiable in $\vx$, and boundary terms vanish under integration by parts. The conditional and marginal paths satisfy \Cref{eq:continuity,eq:conditional_continuity}.
\end{assumption}

% Assumption is boring but needed. i guess.. We assumed no ugly boundary terms and that the densities are nice. Could maybe say this is the usual math hygiene for entropy derivative.

% Extra detail if reviewers ask:
% \begin{equation}
%   \lim_{\|\vx\|\to\infty} p_t(\vx)\bar{\vv}_t(\vx)\log p_t(\vx)\cdot n(\vx)=0,
% \end{equation}
% and same for the conditional density. This is the thing that kills the boundary term.

% \begin{theorem}[Conditional--marginal entropy identity]
% \label{thm:cm_identity}
% Under \Cref{ass:smooth},
% \begin{equation}
%   \frac{\dd}{\dd t}\Ent(Z\mid X_t)
%   =
%   \E_{Z,X_t\mid Z}\!\left[\diver \vv_t(X_t\mid Z)\right]
%   -
%   \E_{X_t}\!\left[\diver \bar{\vv}_t(X_t)\right].
%   \label{eq:cm_identity}
% \end{equation}
% \end{theorem}
\begin{theorem}[Conditional--marginal entropy identity]
    \label{thm:cm_identity}
    Under \Cref{ass:smooth}, \(\displaystyle \frac{\dd}{\dd t}\Ent(Z\mid X_t)=\E_{Z,X_t\mid Z}\!\left[\diver \vv_t(X_t\mid Z)\right]-\E_{X_t}\!\left[\diver \bar{\vv}_t(X_t)\right]\). \inlineeqtag{eq:cm_identity}
    \end{theorem}

% \begin{proof}
% The identity follows from the chain rule for entropy,
% \begin{equation}
%   \Ent(Z\mid X_t)=\Ent(X_t\mid Z)+\Ent(Z)-\Ent(X_t).
% \end{equation}
% The term $\Ent(Z)$ is constant in time. Applying the continuity equation to the marginal density gives
% \begin{equation}
%   \frac{\dd}{\dd t}\Ent(X_t)
%   =
%   \E_{X_t}\!\left[\diver \bar{\vv}_t(X_t)\right].
% \end{equation}
% Applying the same calculation conditionally and then averaging over $Z$ gives
% \begin{equation}
%   \frac{\dd}{\dd t}\Ent(X_t\mid Z)
%   =
%   \E_{Z,X_t\mid Z}\!\left[\diver \vv_t(X_t\mid Z)\right].
% \end{equation}
% Substitution proves the result. \Cref{app:full_derivations} gives the integration-by-parts steps.
% \end{proof}

Use \(\Ent(Z\mid X_t)=\Ent(X_t\mid Z)+\Ent(Z)-\Ent(X_t)\). Since \(\Ent(Z)\) is constant and the continuity equation gives \(\frac{\dd}{\dd t}\Ent(X_t)=\E_{X_t}[\diver\bar{\vv}_t(X_t)]\), with the same calculation conditionally giving \(\frac{\dd}{\dd t}\Ent(X_t\mid Z)=\E_{Z,X_t\mid Z}[\diver\vv_t(X_t\mid Z)]\), substitution yields \Cref{eq:cm_identity}. Full integration-by-parts details are in \Cref{app:full_derivations}.

The theorem is useful because it separates endpoint-specific volume change from population volume change. In a bridge, the conditional field may expand or contract because it is resolving endpoint information, while the marginal field may show a different expansion after averaging over endpoints. The conditional--marginal rate is the difference.

% Maybe here say  about this being the anthagonist of the paper... if we collapse conditional and marginal too early, we kill the bridge signal. The estimator is basically refusing to collapse too early. Sounds too dramatic maybe, keep only if intro needs more story.

\begin{proposition}[Score form of the same rate]
\label{prop:score_form}
Under the assumptions of \Cref{thm:cm_identity}, \\ \(\displaystyle \frac{\dd}{\dd t}\Ent(Z\mid X_t)=-\E_{Z,X_t\mid Z}[(\nabla_{\vx}\log p_t(X_t\mid Z)-\nabla_{\vx}\log p_t(X_t))^\top\vv_t(X_t\mid Z)]\). \inlineeqtag{eq:score_form}
\end{proposition}

\Cref{prop:score_form} is not the estimator used in the high-dimensional experiments, since conditional scores are usually not exposed by pretrained models. It is included because it explains what the divergence contrast measures: the conditional field is weighted by the gap between the conditional score and the marginal score.

\subsection{Gaussian bridges}
\label{sec:gaussian}

For a Gaussian conditional path,
$
  p_t(\vx\mid z)=\mathcal{N}\big(\vx;\mu(z,t),\sigma^2(z,t)I\big),
$
a sample can be written as $X_t=\mu(z,t)+\sigma(z,t)\veps$, with $\veps\sim \mathcal{N}(0,I)$. Differentiating the sample path gives
$
  \vv_t(\vx\mid z)
  =
  \partial_t \mu(z,t)
  +
  \partial_t \sigma(z,t)\frac{\vx-\mu(z,t)}{\sigma(z,t)}.
$
Using the Gaussian score
$
  \nabla_{\vx}\log p_t(\vx\mid z)
  =
  -\frac{\vx-\mu(z,t)}{\sigma^2(z,t)},
$
we obtain the score form
\begin{equation}
  \vv_t(\vx\mid z)
  =
  \partial_t\mu(z,t)
  -
  \sigma(z,t)\partial_t\sigma(z,t)\nabla_{\vx}\log p_t(\vx\mid z).
  \label{eq:gaussian_velocity_score}
\end{equation}

For a Brownian bridge between $x_0$ and $x_1$, write
%\begin{equation}
  $m_t=(1-t)x_0+tx_1$ %\qquad
  and
  $\sigma(t)=\sigma_0\sqrt{t(1-t)}$.
%  \label{eq:brownian_bridge_mean_var}
%\end{equation}
Substitution into \Cref{eq:gaussian_velocity_score} gives the deterministic probability-flow field
\begin{equation}
  \vv_t(\vx\mid x_0,x_1)
  =
  (x_1-x_0)
  +
  \frac{1-2t}{2t(1-t)}(\vx-m_t).
  \label{eq:bb_ode_field}
\end{equation}
The noise scale $\sigma_0$ cancels. This is important because the schedule shape is a property of bridge geometry, not of an arbitrary Brownian scale.

\begin{proposition}[Closed-form conditional divergence]
    \label{prop:bb_divergence}
    For the Brownian-bridge probability-flow field in \Cref{eq:bb_ode_field}, \(\displaystyle \diver \vv_t(\vx\mid x_0,x_1)=d\,\frac{1-2t}{2t(1-t)}\). \inlineeqtag{eq:bb_divergence}
    \end{proposition}

\begin{proof}
The term $x_1-x_0$ is independent of $\vx$ and has zero divergence. The term $m_t$ is also independent of $\vx$. Therefore
\begin{equation}
  \diver \left[\frac{1-2t}{2t(1-t)}(\vx-m_t)\right]
  =
  \frac{1-2t}{2t(1-t)}\diver \vx
  =
  d\,\frac{1-2t}{2t(1-t)}.
\end{equation}
\end{proof}

The magnitude of \Cref{eq:bb_divergence} is singular near both endpoints and vanishes at the midpoint. This is the mathematical source of the U-shaped bridge profile. Because the conditional divergence is spatially constant and endpoint-independent in the Brownian case, averaging over an entropic optimal-transport coupling preserves the same conditional expectation. The full conditional--marginal rate further subtracts the marginal divergence term in \Cref{eq:cm_identity}; the conditional closed form supplies a model-independent anchor and the marginal term adapts it to the learned population path. The U-shape is a theorem for the Gaussian Brownian bridge, not for every learned bridge. Outside this class, the identity remains the general object and the profile must be estimated or derived from the model. Other interpolants can yield flat or monotone analytic profiles, as summarized in \Cref{app:interpolation_profiles}.

\subsection{SDE drift and probability flow}
\label{sec:sde_ode}

For a Brownian bridge, the conditional SDE drift can be written
\begin{equation}
  u^o_t(\vx\mid x_0,x_1)
  =
  (x_1-x_0)
  +
  \frac{1-2t}{t(1-t)}(\vx-m_t),
  \label{eq:bb_sde_drift}
\end{equation}
with conditional score
\begin{equation}
  \nabla_{\vx}\log p_t(\vx\mid x_0,x_1)
  =
  \frac{m_t-\vx}{\sigma_0^2t(1-t)}.
  \label{eq:bb_cond_score}
\end{equation}
The ODE field in \Cref{eq:bb_ode_field} and the SDE drift in \Cref{eq:bb_sde_drift} are distinct conditional objects. They satisfy
\begin{equation}
  u^o_t(\vx\mid x_0,x_1)
  =
  2\vv_t(\vx\mid x_0,x_1)-(x_1-x_0).
  \label{eq:sde_ode_cond_relation}
\end{equation}
The usual probability-flow relation applies to marginal fields after integrating over bridge endpoints. It should not be applied naively to the conditional drift. \Cref{app:sde_ode_details} gives the step-by-step algebra, including the cancellation that explains the factor of two. This distinction is operationally important: using the conditional SDE drift as if it were already the probability-flow field removes the bridge contraction term and makes the conditional profile appear artificially flat.

\section{Estimator, grids, and experimental loops}
\label{sec:estimator}

We organize the method into three layers. The inner layer is mathematical: choose the rate that a bridge actually exposes, namely the conditional divergence minus the marginal divergence. The middle layer is computational: estimate that rate without forming a Jacobian. The outer layer is empirical: use the estimated rate to build a grid, then ask whether the resulting allocation explains the observed behavior of high-dimensional samplers. The layers are tied together because the outer evidence is meaningful only if the computational estimator matches the mathematical rate.

The cond--marg rate is estimated by evaluating both divergence terms in \Cref{eq:cm_identity}. Exact traces are infeasible in high dimension, so the middle loop uses Hutchinson estimation:
\begin{equation}
  \diver \vv_t(\vx)
  =
  \Tr \nabla_{\vx}\vv_t(\vx)
  =
  \E_{\vu}\!\left[\vu^\top\nabla_{\vx}\vv_t(\vx)\vu\right],
  \qquad
  \E[\vu\vu^\top]=I.
  \label{eq:hutch}
\end{equation}
Using the same random probe for the conditional and marginal terms reduces the variance of their difference. After averaging over calibration samples and time points, the rate is smoothed, clipped to a small positive floor, transformed by the default regularized map $\phi(r)=\log(1+r)$, and converted into a grid through \Cref{eq:inverse_cdf}. We also report the raw transform $\phi(r)=r$  as an ablation. The raw grid preserves the uncompressed cond--marg signal, whereas the log transform keeps the ordering of high-rate regions while reducing endpoint domination.

% This should maybe say why same probe matters. If the two traces are close, independent probes make the subtraction noisy for no good reason. Same probe is kind of like paired evaluation. Need not oversell because still variance exists.
% Extra derivation idea:
% \begin{align}
%   \mathrm{Var}[\hat c-\hat m]
%   &=
%   \mathrm{Var}[\hat c]+\mathrm{Var}[\hat m]-2\mathrm{Cov}[\hat c,\hat m].
% \end{align}
% With shared probes the covariance can be positive, so the difference is less noisy. This is an intuition, not a theorem for every field.

\begin{computebox}
\textbf{Computational note.} With $M$ calibration times, $n$ calibration states, and $m$ Hutchinson probes, the estimator uses $O(Mnm)$ derivative-vector products. A typical calibration with 50 times and four probes uses 200 derivative-vector products, which is then amortized across all samples drawn with the same grid. Once the grid is built, sampling at a fixed NFE has the same number of model evaluations as linear, cosine, sigmoid, or power grids. The high-dimensional evidence uses the same inference budgets as the baselines and treats entropy estimation as a calibration step. Memory and arithmetic-intensity details are in \Cref{app:compute}.
\end{computebox}

The output of this construction is a density over time and the grid obtained from its inverse CDF. \Cref{fig:alphaflow_schedule} summarizes the AlphaFlow cond--marg allocation, while \Cref{fig:alphaflow_schedule_geometry} compares the regularized node geometry against linear and against the ten-step DTW alignment.
    \begin{figure}[t]
      \centering
      \includegraphics[height=0.15\textheight, width=0.95\linewidth]{fig/vfinal_AlphaFlow_scheduler_allocation_cond_marg_only.pdf}
      \caption{AlphaFlow cond--marg schedule allocation. The inverse-CDF grid places nodes according to the estimated conditional--marginal rate for 5, 10, and 25 steps; shaded endpoint bands mark the boundary regions used for the schedule-geometry diagnostic. This is grid-allocation evidence, not endpoint-quality evidence.}
      \label{fig:alphaflow_schedule}
    \end{figure}

% The grid is just the normalized information curve turned into steps. Reders should not think it is magic optimizer. It is more like measuring where the path moves information and putting steps there. Need make this sentence less vague if used.

% \begin{figure}[t]
%   \centering
%   \includegraphics[width=0.84\linewidth]{AlphaFlow_cond_marg_u_shape_profile.pdf}
%   \caption{AlphaFlow cond--marg profile. The estimated rate recovers the predicted U-shaped structure and differs sharply from a linear allocation.}
%   \label{fig:alphaflow_profile}
% \end{figure}

    % \begin{figure}[t]
    %   \centering
    %   \includegraphics[width=0.88\linewidth]{AlphaFlow_scheduler_allocation_cond_marg_only.pdf}
    %   \caption{AlphaFlow schedule induced by the cond--marg rate. The profile is converted into node placement by normalized CDF inversion, and BCR summarizes endpoint concentration when visual comparisons are ambiguous.}
    %   \label{fig:alphaflow_schedule}
    % \end{figure}

    % \begin{figure}[t]
    %   \centering
    %   \begin{subfigure}[t]{0.49\linewidth}
    %     \centering
    %     \includegraphics[width=\linewidth]{fig/v3_AlphaFlow_scheduler_allocation_cond_marg_only.pdf}
    %     \caption{Cond--marg allocation.}
    %   \end{subfigure}
    %   \hfill
    %   \begin{subfigure}[t]{0.49\linewidth}
    %     \centering
    %     \includegraphics[width=\linewidth]{fig/v3_alphaflow_log1p_vs_linear_nodes.pdf}
    %     \caption{Cond--marg Log1p regularized versus linear nodes.}
    %   \end{subfigure}
    %   \caption{AlphaFlow cond--marg schedule. Left: inverse-CDF allocation and endpoint BCR. Right: log1p-regularized nodes versus linear at 25 steps. Visual resemblance, but quantitatively different.}
    %   \label{fig:alphaflow_schedule}
    % \end{figure}

Algorithmic summaries are given in \Cref{app:algorithms}. We also define a boundary concentration ratio (BCR) to summarize how much grid mass is placed near the endpoints. Visual inspection of schedule curves can be misleading when two grids look similar but allocate different endpoint density, so BCR gives a quantitative check of boundary concentration. This auxiliary measurement can be useful beyond this work for comparing schedule concentration across models; details and derivation are in \Cref{app:bcr}. We use BCR only as a grid-geometry diagnostic, not as an endpoint-quality metric.

\begin{proposition}[Weighted step-size rule]
    \label{prop:ode_error}
    For an order-$p$ one-step ODE solver with local error constant $C(t)$, the continuous step density minimizing $\int_0^1 C(t)\rho(t)^{-p}\dd t$ subject to $\int_0^1\rho(t)\dd t=1$ is \(\displaystyle \rho^\star(t)\propto C(t)^{1/(p+1)}\). \inlineeqtag{eq:rho_star}
    \end{proposition}

The entropy rate is not asserted to equal $C(t)$ exactly. It is used as a tractable proxy for regions where the vector field compresses, expands, or bends sharply. Therefore, the default entropic grid uses the regularized rate $\log(1+r)$. The raw rate is kept as an ablation because it preserves the mathematical signal but can over-concentrate the grid near singular endpoints. This is the same practical lesson seen across diffusion scheduler and fast-solver work: a principled time signal still needs a numerically stable parametrization at low-NFE \citep{karras2022edm, lin2024common, chen2023importance, lu2022dpm, stancevic2025entropic}. The same rate can matter differently for deterministic and stochastic solvers: ODE integration exposes grid-placement error directly, while SDE noise can smooth some placement errors. The full derivation of \Cref{prop:ode_error} is in \Cref{app:ode_error}.

\section{Results}
\label{sec:evidence}

Before moving to high-dimensional samplers, we ran a controlled two-dimensional pre-flight experiment. We trained residual vector-field models on four transport geometries: continuous--continuous (C-C), continuous--discrete (C-D), discrete--continuous (D-C), and discrete--discrete (D-D), across bridge scales. We then built grids from the measured entropy-rate profiles. These learned probes recover the boundary-heavy Brownian-bridge pattern, place more nodes near endpoints, and show the largest gains at low NFE before schedules converge. We report both ODE and SDE sweeps, but use the probability-flow ODE for the high-dimensional study because it removes sampling noise and matches the cond--marg estimator. Full transport visuals and ODE/SDE curves are in \Cref{app:figures} and the controlled transport setup and entropy maps are summarized in in  \Cref{fig:controlled_transport_setup_appendix,fig:mmd_steps_appendix}.

% Although these models are small, their fields are learned rather than closed forms and, therefore, they test whether the boundary-heavy Brownian-bridge pattern still appears after training. We show that the measured rates are boundary-heavy, the inverse-CDF grids place more nodes near the endpoints, and the low-NFE gains are largest before schedule differences shrink at higher step counts. We also compared ODE and SDE sampling in this synthetic setting. Both solvers remain schedule-sensitive; however, the probability-flow ODE is the clearer interface for the later high-dimensional study because it removes sampling noise and matches the deterministic field used by the cond--marg estimator. The toy and synthetic appendix in \Cref{app:figures} collects the corresponding transport setup, entropy maps, and ODE/SDE step curves in \Cref{fig:mmd_steps_appendix}; the controlled transport setup and entropy maps are summarized in in  \Cref{fig:controlled_transport_setup_appendix,fig:mmd_steps_appendix}.

\begin{table}[!t]
  \vspace{-0.75em}
  \centering
  \small
  \setlength{\tabcolsep}{3pt}
  \caption{Synthetic low-NFE results. Left: learned ODE-Heun MMD averaged across steps, scenarios, $\sigma$, and seeds. Right: entropic-vs-linear improvement with bootstrap intervals.}
  \label{tab:low_nfe_improvement}

  \begin{minipage}[t]{0.49\linewidth}
    \centering
    \scriptsize
    \renewcommand{\arraystretch}{1.05}
    \begin{tabular*}{\linewidth}{@{\extracolsep{\fill}}lcc@{}}
      \toprule
      Scheduler & MMD $\downarrow$ ($\times 10^{-3}$) & Entropic vs (\%) \\
      \midrule
      Linear  & 6.028$\pm$0.603 & 8.1 [6.6, 9.7] \\
      Power-2 & 6.620$\pm$1.070 & 16.4 [14.1, 18.6] \\
      Power-3 & 6.912$\pm$1.216 & 19.9 [17.6, 22.3] \\
      Log     & 8.164$\pm$1.855 & 32.2 [29.7, 34.6] \\
      Cosine  & 5.548$\pm$0.152 & 0.2 [0.1, 0.3] \\
      Sigmoid & 5.639$\pm$0.339 & 1.8 [0.5, 3.1] \\
      \begingroup\setlength{\fboxsep}{1.5pt}\colorbox{yellow!30}{\textbf{Entropic}}\endgroup
      &
      \begingroup\setlength{\fboxsep}{1.5pt}\colorbox{yellow!30}{\textbf{5.537$\pm$0.131}}\endgroup
      &
      \begingroup\setlength{\fboxsep}{1.5pt}\colorbox{yellow!30}{\textbf{---}}\endgroup \\
      \bottomrule
    \end{tabular*}
  \end{minipage}
  \hfill
  \begin{minipage}[t]{0.47\linewidth}
    \centering
    \scriptsize
    \renewcommand{\arraystretch}{1.05}
    \begin{tabular*}{\linewidth}{@{\extracolsep{\fill}}ccc@{}}
      \toprule
      Steps & ODE-Heun (\%) & SDE-Heun (\%) \\
      \midrule
      10 & 18.1 [-4.9, 35.9]  & 22.7 [7.3, 34.9] \\
      25 & 12.3 [-10.8, 31.1] & 13.1 [-3.7, 27.0] \\
      50 & 7.0 [-16.8, 26.7]  & 10.7 [-5.2, 24.8] \\
      \bottomrule
    \end{tabular*}
  \end{minipage}
  \vspace{-0.75em}
\end{table}
% \begin{table}[t]
% \centering
% \small
% \caption{Step-averaged MMD for learned ODE-Heun across $n_{\mathrm{steps}} \in \{10,25,50,100,200\}$, aggregated over scenarios, $\sigma$, and seeds. MMD is scaled by $10^{-3}$; the entropic-vs column reports MMD improvement relative to each scheduler.}
% \label{tab:toy-ode-heun-mmd}
%   \begin{tabular*}{\linewidth}{@{\extracolsep{\fill}}lcc@{}}
%     \toprule
%     Scheduler & MMD $\downarrow$ ($\times 10^{-3}$) & Entropic vs (\%) \\
%     \midrule
%     Linear            & 6.028$\pm$0.603          & 8.1 [6.6, 9.7] \\
%     Power-2           & 6.620$\pm$1.070          & 16.4 [14.1, 18.6] \\
%     Power-3           & 6.912$\pm$1.216          & 19.9 [17.6, 22.3] \\
%     Log               & 8.164$\pm$1.855          & 32.2 [29.7, 34.6] \\
%     Cosine            & 5.548$\pm$0.152          & 0.2 [0.1, 0.3] \\
%     Sigmoid           & 5.639$\pm$0.339          & 1.8 [0.5, 3.1] \\
%     \textbf{Entropic} & \textbf{5.537$\pm$0.131} & --- \\
%     \bottomrule
%     \end{tabular*}
% \end{table}

% \begin{table}[H]
%     \caption{Entropic vs linear MMD improvement at low NFE (supports H3 and H5). Values are mean \% with bootstrap interval across scenarios, model sizes, sigmas, and seeds.}
%     \label{tab:low_nfe_improvement}
%     \centering
%     \begin{small}
%         \begin{tabular*}{\linewidth}{@{\extracolsep{\fill}}ccc@{}}
%             \toprule
%             Steps & ODE-Heun (\% improvement) & SDE-Heun (\% improvement) \\
%             \midrule
%             10    & 18.1 [-4.9, 35.9]         & 22.7 [7.3, 34.9]          \\
%             25    & 12.3 [-10.8, 31.1]        & 13.1 [-3.7, 27.0]         \\
%             50    & 7.0 [-16.8, 26.7]         & 10.7 [-5.2, 24.8]         \\
%             \bottomrule
%         \end{tabular*}
%     \end{small}
% \end{table}

% We next test the same allocation signal in two high-dimensional settings. EDM/CIFAR-10 has $d=32\times32\times3=3072$ coordinates. Although EDM is not a bridge model and is not trained on the entropy objective, it is a useful \textit{mismatch case}: it tests whether a cond--marg-inspired grid can still help when the theory is not exactly matched \citep{karras2022edm}. AlphaFlow operates on protein ensembles with more than $5\times 10^4$ effective coordinates -- across sequence length, atom representation, and ensemble size -- in the evaluated structures. It is closer to the bridge setting because its flow-matching sampler lets us estimate the cond--marg geometry targeted by the theory, although it was not trained with the entropy criterion \citep{jing2024alphafoldmeetsflowmatching}. The high-dimensional analysis focuses on ODE sampling: it is the inference exposed by AlphaFlow and it matches the estimator. We evaluate the protein evidence on CAMEO22 and ATLAS targets \citep{ha2015cameo, stam2024atlas}. The evaluation funnel in \Cref{fig:alphaflow_testing_sankey_appendix} shows 120 available proteins across the two benchmark sources, of which 115 were tested after filtering. More details are in \Cref{app:plddt_metric}.
We next test the same allocation signal in two high-dimensional settings. EDM/CIFAR-10 is a mismatch case: it is not a bridge model and was not trained with the entropy objective, but it tests whether a tempered cond--marg-inspired grid can help a pretrained image sampler \citep{karras2022edm}. AlphaFlow operates on protein ensembles and is closer to the bridge setting because its flow-matching sampler exposes the fields needed for the cond--marg estimator, though it was also not trained with this criterion \citep{jing2024alphafoldmeetsflowmatching}. We use ODE sampling throughout, evaluate CAMEO22 and ATLAS proteins\citep{ha2015cameo, stam2024atlas}, and report the evaluation funnel in \Cref{app:plddt_metric}.



% The main evidence below tests whether that story survives high-dimensional settings. EDM operates on CIFAR-10 images with $d=32\times32\times3=3072$ coordinates and is a diffusion-style image model outside the bridge assumptions. AlphaFlow operates on protein ensembles with more than $5\times 10^4$ effective coordinates in the evaluated structures and is closer to a bridge-like flow, but it was trained with its original objective rather than the entropy objective studied here.




% \begin{computebox}
% \textcolor{red}{\textbf{Update note.} Additional postprocessing jobs are done. Check the website \href{https://www.trentini.fyi/entropic}{https://www.trentini.fyi/entropic} for now while I update the manuscript.}
% \end{computebox}

\textbf{AlphaFlow.} The cond--marg profile in \Cref{fig:entropy_curves} is strongly U-shaped: the intervals $[0,0.1]$ and $[0.9,1]$ each carry about $31.1\%$ of the normalized mass, while the center window $[0.4,0.6]$ carries about $5.0\%$.    The resulting grid allocation is shown in \Cref{fig:alphaflow_schedule}; the log1p node geometry and ten-step DTW alignment are shown in \Cref{fig:alphaflow_schedule_geometry}. Protein structures lie on a thin constrained manifold, where bond geometry, steric exclusions, chirality, and fold commitment can make both endpoints stiff. This motivates boundary allocation, but it does not by itself prove endpoint-quality improvement. Endpoint predicted local distance difference test (pLDDT) is mixed in the available AlphaFlow sweeps summarized in \Cref{tab:alphaflow-plddt-ode}: sigmoid and cosine are strongest in many displayed cells, raw entropic scheduling is strongest in the medium and large five-step cells, and raw entropic pLDDT degrades sharply at 10 and 25 steps. pLDDT is an AlphaFold-style per-residue local distance difference test, and we use it here as an endpoint confidence proxy rather than as a measure of schedule geometry \citep{jumper2021,mariani2013lddt}. Our method leads pLDDT performance on both medium and large-sized proteins ($97.62\pm0.22$ and $97.91\pm0.14$) at the low 5-step NFE count, and it is close to the sigmoid leader for small proteins ($95.06\pm1.12$ versus $95.26\pm0.99$). Its limitations in this paper are summarized in \Cref{app:plddt_metric}.  Additional AlphaFlow grid and structure-evolution visualizations are kept in the Evaluation Grids appendix, including the 25-step AlphaFlow trajectory diagnostic in \Cref{app:evaluation_grids}; \Cref{fig:alphaflow_evolution_25steps_appendix} explains how endpoint plDDT should be interpreted so that we keep our emphasis on the schedule geometry.

% Power grids give the strongest endpoint predicted local distance difference test (plDDT) in the available AlphaFlow sweeps, while entropic grids remain close to cosine and improve several paired diagnostics over linear.

% The paired entropic advantages over linear are $0.181$, $0.225$, and $0.278$ pLDDT at 5, 10, and 25 steps, with positive bootstrap intervals. Diversity advantages over linear are $0.074$, $0.058$, and $0.059$. RMSD-to-final timing diagnostics also favor the cond--marg schedule over several baselines, with advantages against power grids of $0.161$, $0.207$, and $0.259$ at 5, 10, and 25 steps.



% \begin{table}[t]
% \centering
% \scriptsize
% \setlength{\tabcolsep}{2.2pt}
% \caption{AlphaFlow pLDDT scores under ODE-Heun sampling. Higher is better. Entries report mean pLDDT with 95\% confidence intervals. This is endpoint-confidence evidence, not direct schedule-geometry evidence. Yellow marks the best displayed scheduler within each size and NFE column; bold marks the reported entropic scheduler.}
% \label{tab:alphaflow-plddt-ode}
% \begin{tabular*}{\linewidth}{@{\extracolsep{\fill}}lccccccccc@{}}
%     \toprule
%     & \multicolumn{3}{c}{Small} & \multicolumn{3}{c}{Medium} & \multicolumn{3}{c}{Large} \\
%     \cmidrule(lr){2-4}\cmidrule(lr){5-7}\cmidrule(l){8-10}
%     Scheduler & 5 & 10 & 25 & 5 & 10 & 25 & 5 & 10 & 25 \\
%     \midrule
%     Linear  & \alphaplddtcell{89.97}{1.83} & \alphaplddtcell{94.15}{0.88} & \alphaplddtcell{95.82}{0.51} & \alphaplddtcell{93.10}{0.61} & \alphaplddtcell{95.08}{0.31} & \alphaplddtcell{96.25}{0.12} & \alphaplddtcell{94.27}{0.57} & \alphaplddtcell{94.30}{0.30} & \alphaplddtcell{94.12}{0.12} \\
%     Cosine  & \alphaplddtcell{92.73}{1.44} & \alphaplddtcell{96.08}{0.66} & \alphabestplddtcell{96.07}{0.54} & \alphaplddtcell{95.23}{0.47} & \alphaplddtcell{97.57}{0.13} & \alphaplddtcell{97.54}{0.06} & \alphaplddtcell{95.92}{0.45} & \alphaplddtcell{97.09}{0.13} & \alphaplddtcell{96.39}{0.07} \\
%     Sigmoid & \alphabestplddtcell{95.26}{0.99} & \alphabestplddtcell{96.15}{0.80} & \alphaplddtcell{95.64}{0.71} & \alphaplddtcell{97.49}{0.26} & \alphabestplddtcell{98.11}{0.08} & \alphabestplddtcell{97.55}{0.06} & \alphaplddtcell{97.82}{0.23} & \alphabestplddtcell{97.95}{0.08} & \alphabestplddtcell{96.67}{0.06} \\
%     Power-2 & \alphaplddtcell{84.81}{2.64} & \alphaplddtcell{90.44}{1.82} & \alphaplddtcell{94.87}{0.78} & \alphaplddtcell{90.92}{0.80} & \alphaplddtcell{93.99}{0.52} & \alphaplddtcell{95.27}{0.22} & \alphaplddtcell{92.54}{0.73} & \alphaplddtcell{94.89}{0.45} & \alphaplddtcell{93.52}{0.20} \\
%     \textbf{Entropic} & \alphaplddtcellbold{95.06}{1.12} & \alphaplddtcellbold{78.79}{3.42} & \alphaplddtcellbold{75.45}{3.81} & \alphabestplddtcellbold{97.62}{0.22} & \alphaplddtcellbold{87.85}{1.11} & \alphaplddtcellbold{86.02}{1.25} & \alphabestplddtcellbold{97.91}{0.14} & \alphaplddtcellbold{90.26}{0.68} & \alphaplddtcellbold{89.26}{0.76} \\
%     \bottomrule
% \end{tabular*}
% \end{table}


% Important: do not say AlphaFlow is won by entropic. It is more complicated. The nice thing is the profile and allocation, but pLDDT endpoint likes power. This is not a failure if we explain that training objective and metric are different thing than grid geometry.

% Maybe add later..."the model was not trained to make the entropic grid optimal", but not too defensive.

% \begin{table}[t]
%   \caption{AlphaFlow endpoint pLDDT on pooled high-dimensional protein evaluations. Higher is better. Power grids are strongest on endpoint confidence, while entropy gives bridge-aligned allocation diagnostics and remains close to the main baselines.}
%   \label{tab:alphaflow}
%   \centering
%   \small
%   \begin{tabular}{lrrr}
%     \toprule
%     Scheduler & 5 steps & 10 steps & 25 steps \\
%     \midrule
%     Linear & 86.1 & 88.1 & 89.0 \\
%     Sigmoid & 86.9 & 88.8 & 89.7 \\
%     Power 2 & \textbf{88.6} & \textbf{90.3} & \textbf{91.0} \\
%     Entropic log1p & 86.3 & 88.3 & 89.3 \\
%     \bottomrule
%   \end{tabular}
%   \vspace{0.35em}
%   {\footnotesize\textcolor{red}{\textbf{TODO $\star$} Waiting for final jobs to add confidence intervals. Nice-to-have, non-essential metrics include divergence and tau-curve diagnostics, FID/sample/timing summaries, protein endpoint quality metrics, structural validity metrics, cross-scheduler evidence, and aggregated confidence-interval tables.}}
% \end{table}

% \begin{figure}[t]
%   \centering
%   \includegraphics[width=0.78\linewidth]{AlphaFlow_entropic_cond_marg_protein_evolution_10steps_3d_7lp1_A.pdf}
%   \caption{Protein trajectory diagnostic for the ten-step conditional--marginal grid. This rendering is used to inspect the trajectory and grid geometry, not as proof of independent endpoint-quality improvement.}
%   \label{fig:protein_static_main}
% \end{figure}

\begin{table}[H]
\centering
\scriptsize
\setlength{\tabcolsep}{2.2pt}
\caption{AlphaFlow endpoint pLDDT under ODE-Heun. Values are means with 95\% CIs; higher is better. Highlighting marks the best scheduler per column and bold marks entropic.}
\label{tab:alphaflow-plddt-ode}
\begin{tabular*}{\linewidth}{@{\extracolsep{\fill}}lccccccccc@{}}
    \toprule
    & \multicolumn{3}{c}{\begin{tabular}[c]{@{}c@{}}Small \,\,{\scriptsize $\le 50$ aa}\end{tabular}}
    & \multicolumn{3}{c}{\begin{tabular}[c]{@{}c@{}}Medium \,\, {\scriptsize 51--400 aa}\end{tabular}}
    & \multicolumn{3}{c}{\begin{tabular}[c]{@{}c@{}}Large \,\,{\scriptsize $>400$ aa}\end{tabular}} \\
    \cmidrule(lr){2-4}\cmidrule(lr){5-7}\cmidrule(l){8-10}
    Scheduler & 5 & 10 & 25 & 5 & 10 & 25 & 5 & 10 & 25 \\
    \midrule
    Linear  & \alphaplddtcell{89.97}{1.83} & \alphaplddtcell{94.15}{0.88} & \alphaplddtcell{95.82}{0.51} & \alphaplddtcell{93.10}{0.61} & \alphaplddtcell{95.08}{0.31} & \alphaplddtcell{96.25}{0.12} & \alphaplddtcell{94.27}{0.57} & \alphaplddtcell{94.30}{0.30} & \alphaplddtcell{94.12}{0.12} \\
    Cosine  & \alphaplddtcell{92.73}{1.44} & \alphaplddtcell{96.08}{0.66} & \alphabestplddtcell{96.07}{0.54} & \alphaplddtcell{95.23}{0.47} & \alphaplddtcell{97.57}{0.13} & \alphaplddtcell{97.54}{0.06} & \alphaplddtcell{95.92}{0.45} & \alphaplddtcell{97.09}{0.13} & \alphaplddtcell{96.39}{0.07} \\
    Sigmoid & \alphabestplddtcell{95.26}{0.99} & \alphabestplddtcell{96.15}{0.80} & \alphaplddtcell{95.64}{0.71} & \alphaplddtcell{97.49}{0.26} & \alphabestplddtcell{98.11}{0.08} & \alphabestplddtcell{97.55}{0.06} & \alphaplddtcell{97.82}{0.23} & \alphabestplddtcell{97.95}{0.08} & \alphabestplddtcell{96.67}{0.06} \\
    Power-2 & \alphaplddtcell{84.81}{2.64} & \alphaplddtcell{90.44}{1.82} & \alphaplddtcell{94.87}{0.78} & \alphaplddtcell{90.92}{0.80} & \alphaplddtcell{93.99}{0.52} & \alphaplddtcell{95.27}{0.22} & \alphaplddtcell{92.54}{0.73} & \alphaplddtcell{94.89}{0.45} & \alphaplddtcell{93.52}{0.20} \\
    \textbf{Entropic} & \alphaplddtcellbold{95.06}{1.12} & \alphaplddtcellbold{78.79}{3.42} & \alphaplddtcellbold{75.45}{3.81} & \alphabestplddtcellbold{97.62}{0.22} & \alphaplddtcellbold{87.85}{1.11} & \alphaplddtcellbold{86.02}{1.25} & \alphabestplddtcellbold{97.91}{0.14} & \alphaplddtcellbold{90.26}{0.68} & \alphaplddtcellbold{89.26}{0.76} \\
    \bottomrule
\end{tabular*}
\end{table}

\textbf{EDM.} EDM is a mismatch test because it lies outside the Schr\"odinger-bridge setting \citep{karras2022edm}. We use log-tempered cond--marg entropy as the default scheduler and raw entropy as the ablation: raw allocation over-concentrates endpoints and performs poorly, while \(\log(1+r)\) keeps the high-rate regions without exhausting the low-NFE budget at singular endpoints. \Cref{tab:edm-fid-ode} reports the full ODE-Heun FID sweep: at five steps entropic achieves $186.26\pm3.97$, compared with $200.52\pm2.91$ for linear, $238.03\pm5.29$ for cosine, and $355.06\pm2.20$ for sigmoid, while 10/25-step schedules are close.     Additional trajectory-grid diagnostics are in \Cref{app:evaluation_grids}. The schedule-geometry comparison in \Cref{fig:alphaflow_schedule_geometry} shows the same calibration point: log1p stays close to linear but remains endpoint-biased, whereas raw cond--marg allocation requires stronger warping.

% \textbf{EDM.} EDM lies outside the Schr\"odinger-bridge setting, so it is the useful counterweight to AlphaFlow \citep{karras2022edm}. We treat log-tempered cond--marg entropy as the default scheduler and the raw version as the ablation. The ablation is useful because it shows what happens when the bridge rate is used without regularization: raw entropy over-concentrates the grid near the endpoints, leaves too few steps for the rest of the trajectory, and performs poorly. This drop is consistent with prior scheduler work showing that low-NFE diffusion and ODE samplers are sensitive to time parametrization, noise schedules, and endpoint concentration \citep{karras2022edm,lin2024common,chen2023importance,lu2022dpm,lu2022dpmsolverpp}. Log tempering regularizes this effect. It compresses the sharp cond--marg U-shape rather than discarding it. \Cref{tab:edm-fid-ode} lists the full ODE-Heun Inception FID sweep; at five steps entropic log1p achieves FID $186.26\pm3.97$, compared with $200.52\pm2.91$ for linear, $238.03\pm5.29$ for cosine, and $355.06\pm2.20$ for sigmoid. At 10 and 25 steps the best schedules are close. The trajectory comparison in \Cref{fig:edm_images} is kept in the Evaluation Grids appendix as \Cref{fig:edm_images}; it shows the five-step regime where this difference is most visible. Visually, the $\log(1+r)$ grid is close to linear. Quantitatively, however, it is an endpoint-biased version of linear whereas the raw cond--marg grid is strongly endpoint-concentrated. For the five-step grids used in the figure, the mean distance from linear is $0.0169$ for $\log(1+r)$ and $0.0969$ for raw cond--marg allocation; the interval-width coefficient of variation is $0.143$ for $\log(1+r)$ and $0.935$ for raw allocation. Hence, $\log(1+r)$ is not linear, but it is regularized enough to avoid wasting the low-NFE budget at the singular endpoints. Additional EDM grid visualizations are in \Cref{app:evaluation_grids}; auxiliary heatmaps are omitted from the appendix because the main claim is already captured by \Cref{tab:edm-fid-ode} and the trajectory grids. The appendic DTW diagnostic in \Cref{fig:schedule_dtw_warping_paths_appendix} quantifies the same point: at 10 steps, cond--marg $\log(1+r)$ stays close to linear but remains distinct, whereas raw cond--marg allocation requires stronger index warping.
% The diagnostic available in \Cref{fig:schedule_dtw_warping_paths_appendix} adds a 10-step dynamic-time-warping (DTW) diagnostic for AlphaFlow's scheduler-path geometry. This diagnostic makes the scheduler-similarity claim more precise: cond--marg $\log(1+r)$ stays close to the linear diagonal but remains a distinct time series, whereas raw cond--marg allocation requires stronger index warping. For the AlphaFlow 10-step grids the normalized DTW distance to linear is $0.0227$ for cond--marg $\log(1+r)$ and $0.0525$ for raw cond--marg allocation \footnote{The appendix figure should be read as schedule-geometry evidence, not as an endpoint-quality metric.}.

% \textbf{EDM.} EDM lies outside the Schr\"odinger-bridge setting, so it is the useful counterweight to AlphaFlow \citep{karras2022edm}. If entropy alone were the whole story, raw entropy would be the natural winner. The tests reject that premise. Raw entropy over-concentrates the grid and performs poorly. Log tempering is the stable variant. At five ODE-Heun steps it achieves FID $184.9\pm4.8$, compared with $198.4\pm1.8$ for linear, $234.3\pm2.7$ for cosine, and $354.4\pm2.3$ for sigmoid. At 10 and 25 steps the best schedules are close. The image grid in \Cref{fig:edm_images} shows the five-step regime where this difference is most visible. Additional EDM grid visualizations are in \Cref{app:evaluation_grids}; auxiliary heatmaps are omitted from the appendix because the main claim is already captured by the FID table and trajectory grids.


% EDM is the stress test that almost breaks the simple story. Raw entropy is too extreme. log1p is like admitting the math signal needs numerical calibration. This PROBABLY GOOD because it prevents over claiming!

% Extra possible equation:
% \begin{equation}
%   r_{\mathrm{log}}(t)=\log(1+r_{\mathrm{raw}}(t))
% \end{equation}
% This preserves ordering but compresses endpoint peaks.

% \begin{table}[t]
%   \caption{EDM image generation at low NFE. Lower FID is better. Log-tempered entropy is strongest at five ODE steps and competitive as NFE increases.}
%   \label{tab:edm}
%   \centering
%   \small
%   \begin{tabular}{llrrr}
%     \toprule
%     Solver & Scheduler & 5 steps & 10 steps & 25 steps \\
%     \midrule
%     ODE Heun & Linear & $198.4\pm1.8$ & $173.1\pm3.6$ & $171.3\pm4.6$ \\
%     ODE Heun & Cosine & $234.3\pm2.7$ & $177.2\pm2.9$ & $\mathbf{170.9\pm4.4}$ \\
%     ODE Heun & Sigmoid & $354.4\pm2.3$ & $207.9\pm3.9$ & $173.4\pm2.9$ \\
%     ODE Heun & Entropic raw & $323.0\pm6.3$ & $290.2\pm6.0$ & $227.8\pm1.4$ \\
%     ODE Heun & Entropic log1p & $\mathbf{184.9\pm4.8}$ & $\mathbf{172.4\pm4.9}$ & $170.9\pm4.5$ \\
%     \midrule
%     SDE Heun & Linear & $\mathbf{382.2\pm4.1}$ & $356.1\pm5.4$ & $\mathbf{336.9\pm7.5}$ \\
%     SDE Heun & Entropic log1p & $387.9\pm3.6$ & $\mathbf{355.4\pm5.0}$ & $336.9\pm6.8$ \\
%     \bottomrule
%   \end{tabular}
% \end{table}



% \begin{table}[t]
%   \caption{EDM image-generation FID scores under ODE-Heun sampling. Lower is better. Entries report mean Inception FID with 95\% confidence intervals over five seeds. The entropic default is the cond--marg log1p scheduler; raw entropy is kept as an ablation because it over-concentrates the grid near the endpoints. Yellow marks the best displayed scheduler for each NFE regime.}
%   \label{tab:edm-fid-ode}
%   \centering
%   \small
%   \begin{tabular*}{\linewidth}{@{\extracolsep{\fill}}lccc@{}}
%     \toprule
%     Scheduler & 5 steps & 10 steps & 25 steps \\
%     \midrule
%     ODE + Linear  & \edmfidcell{200.52}{2.91} & \edmfidcell{172.78}{2.83} & \edmfidcell{170.85}{3.50} \\
%     ODE + Cosine  & \edmfidcell{238.03}{5.29} & \edmfidcell{176.58}{2.49} & \edmfidcell{170.80}{3.36} \\
%     ODE + Sigmoid & \edmfidcell{355.06}{2.20} & \edmfidcell{212.15}{8.03} & \edmfidcell{172.68}{2.39} \\
%     ODE + Power-2 & \edmfidcell{316.54}{2.68} & \edmfidcell{189.13}{3.51} & \edmfidcell{171.28}{4.06} \\
%     ODE + Entropic (raw) & \edmfidcell{322.58}{4.80} & \edmfidcell{290.11}{4.94} & \edmfidcell{227.96}{1.43} \\
%     \textbf{ODE + Entropic} & \edmbestfidcell{186.26}{3.97} & \edmbestfidcell{172.10}{3.74} & \edmbestfidcell{170.78}{3.42} \\
%     \bottomrule
%   \end{tabular*}
% \end{table}
% \begin{figure}[H]
%   \centering
%   \includegraphics[width=\linewidth]{alphaflow_schedule_dtw_warping_paths_10step.pdf}
%   \caption{AlphaFlow 10-NFE dynamic-time-warping alignment paths between scheduler time series. Columns are query schedulers; rows use raw cond--marg and log-tempered cond--marg schedules as AlphaFlow references. The dark path is the optimal DTW alignment, and the color scale gives local index-pair distance.}
%   \label{fig:schedule_dtw_warping_paths_appendix}
% \end{figure}
\begin{table}[!t]
  \caption{EEDM/CIFAR-10 ODE-Heun FID. Values are means with 95\% CIs over five seeds; lower is better. Entropic uses cond--marg log1p, with raw entropy shown as an ablation.}
  \label{tab:edm-fid-ode}
  \centering
  \scriptsize
  \setlength{\tabcolsep}{1.8pt}
  \begin{adjustbox}{max width=\linewidth}
  \begin{tabular}{@{}lccc@{}}
    \toprule
    Scheduler & 5 steps & 10 steps & 25 steps \\
    \midrule
    ODE + Linear  & \edmfidcell{200.52}{2.91} & \edmfidcell{172.78}{2.83} & \edmfidcell{170.85}{3.50} \\
    ODE + Cosine  & \edmfidcell{238.03}{5.29} & \edmfidcell{176.58}{2.49} & \edmfidcell{170.80}{3.36} \\
    ODE + Sigmoid & \edmfidcell{355.06}{2.20} & \edmfidcell{212.15}{8.03} & \edmfidcell{172.68}{2.39} \\
    ODE + Power-2 & \edmfidcell{316.54}{2.68} & \edmfidcell{189.13}{3.51} & \edmfidcell{171.28}{4.06} \\
    ODE + Entropic (raw) & \edmfidcell{322.58}{4.80} & \edmfidcell{290.11}{4.94} & \edmfidcell{227.96}{1.43} \\
    \textbf{ODE + Entropic} & \edmbestfidcell{186.26}{3.97} & \edmbestfidcell{1å72.10}{3.74} & \edmbestfidcell{170.78}{3.42} \\
    \bottomrule
  \end{tabular}
  \end{adjustbox}
\end{table}

  \begin{figure*}[t]
    \centering
    \includegraphics[width=\textwidth]{fig/vfinal_alphaflow_schedule_geometry_lanes.pdf}
    \caption{AlphaFlow schedule-geometry diagnostics. Left: the conditional--marginal log1p grid remains \textit{visually} close to the linear grid across 5, 10, and 25 steps while keeping
  endpoint bias quantitatively. Right: dynamic-time-warping paths at 10 steps quantify this distinction; the log1p grid has lower distance to linear than the raw cond--marg grid, while raw cond--marg requires stronger index warping.}
    \label{fig:alphaflow_schedule_geometry}
  \end{figure*}



    \textbf{Cross-domain allocation.} \Cref{fig:entropy_curves} puts the EDM and AlphaFlow entropy profiles on the same page, \Cref{fig:alphaflow_schedule} shows the AlphaFlow allocation, and \Cref{fig:alphaflow_schedule_geometry} shows the log1p-vs-linear node geometry together with the ten-step DTW alignment. In these diagnostics the bridge-aware quantity is the two-term rate
\begin{equation}
  r_{\cm}(t)
  =
  \left|
  \E_{Z,X_t\mid Z}\!\left[\diver\vv_t(X_t\mid Z)\right]
  -
  \E_{X_t}\!\left[\diver\bar{\vv}_t(X_t)\right]
  \right|.
  \label{eq:results_cm_rate}
\end{equation}
The profiles are not the same, which is the point of measuring the entropy rate. EDM needs tempering because its raw profile is too sharp for stable low-NFE integration. AlphaFlow shows a bridge-like U-shape, but endpoint confidence remains influenced by the training objective and protein-specific geometry. The evidence separates three ideas that are easy to conflate: the entropy profile, the numerical grid derived from that profile, and the downstream metric used to judge the generated sample. Additional grid and evolution visualizations for EDM and AlphaFlow are in \Cref{app:evaluation_grids}.

% \begin{figure}[t]
%   \centering
%   \begin{subfigure}[t]{0.49\linewidth}
%     \centering
%     \includegraphics[width=\linewidth]{figures__EDM_entropy_curve_shape.pdf}
%     \caption{EDM entropy curve.}
%   \end{subfigure}
%   \hfill
%   \begin{subfigure}[t]{0.49\linewidth}
%     \centering
%     \includegraphics[width=\linewidth]{AlphaFlow_cond_marg_u_shape_profile.pdf}
%     \caption{AlphaFlow conditional--marginal curve.}
%   \end{subfigure}
%   \caption{Entropy profiles for the two high-dimensional stress tests. EDM is a mismatch case that requires tempering; AlphaFlow shows the bridge-like U-shape predicted by the Brownian calculation.}
%   \label{fig:entropy_curves}
% \end{figure}

    % \begin{figure}[t]
    %   \centering
    %   \begin{subfigure}[t]{0.49\linewidth}
    %     \centering
    %     \includegraphics[width=\linewidth]{figures__EDM_entropy_curve_shape.pdf}
    %     \caption{EDM entropy curve.}
    %   \end{subfigure}
    %   \hfill
    %   \begin{subfigure}[t]{0.49\linewidth}
    %     \centering
    %     \includegraphics[width=\ladded_analysis__EDM_scheduler_allocation_cond_marg_onlyinewidth]{AlphaFlow_cond_marg_u_shape_profile_density.pdf}
    %     \caption{AlphaFlow cond--marg density.}
    %   \end{subfigure}
    %   \caption{Entropy-profile diagnostics for the two high-dimensional stress tests. Left: EDM has a sharp profile that requires log tempering for stable low-NFE sampling. Right: AlphaFlow shows the bridge-like boundary-heavy conditional--marginal density predicted by the Brownian-bridge calculation.}
    %   \label{fig:entropy_curves}
    % \end{figure}

    % \begin{figure}[t]
    %   \centering
    %   \includegraphics[width=\linewidth]{fig/edm_alphaflow_entropy_profiles_tarchic.png}
    %   \caption{Entropy-profile diagnostics for the two high-dimensional stress tests. Left: EDM has a sharp profile that requires log tempering for stable low-NFE sampling. Right: AlphaFlow shows the bridge-like boundary-heavy conditional--marginal density predicted by the Brownian-bridge calculation. Shaded endpoint bands mark the BCR regions used to summarize endpoint concentration.}
    %   \label{fig:entropy_curves}
    % \end{figure}

% \begin{figure}[t]
%   \centering
%   % \includegraphics[width=0.88\linewidth]{added_analysis__EDM_scheduler_allocation_cond_marg_only.pdf}
%   % \vspace{0.6em}
%   \includegraphics[width=0.88\linewidth]{AlphaFlow_scheduler_allocation_cond_marg_only.pdf}
%   \caption{AlphaFlow cond--marg. The cond--marg rate is convereted into node placement by normalized CDF inversion, BCR -- shaded regions -- is used as a scalar summary of endpoint concentration when visual comparisons are ambiguous.}
%   \label{fig:allocation}
% \end{figure}

% \begin{figure}[t]
%   \centering
%   \begin{subfigure}[t]{0.49\linewidth}
%     \centering
%     \includegraphics[width=\linewidth]{cond_marg_log1p_vs_linear_edm_alphaflow.pdf}
%     \caption{Log1p versus linear.}
%   \end{subfigure}
%   \hfill
%   \begin{subfigure}[t]{0.49\linewidth}
%     \centering
%     \includegraphics[width=\linewidth]{cond_marg_timeseries_10step_all_schedulers_edm_alphaflow.pdf}
%     \caption{Ten-step scheduler time series.}
%   \end{subfigure}
%   \caption{Auxiliary scheduler diagnostics for EDM and AlphaFlow. These panels make the same point as BCR: visually similar grids can have different endpoint concentration and different time-series geometry.}
%   \label{fig:cross_domain_scheduler_diagnostics}
% \end{figure}
    % \begin{figure}[t]
    %   \centering
    %   \includegraphics[width=0.72\linewidth]{alphaflow_log1p_vs_linear_nodes.pdf}
    %   \caption{AlphaFlow Log1p versus linear scheduler diagnostic. The panel compares the linear grid with the AlphaFlow conditional--marginal log1p grid while keeping endpoint-quality metrics separate from schedule geometry.}
    %   \label{fig:cross_domain_scheduler_diagnostics}
    % \end{figure}




\section{Discussion}
\label{sec:discussion}

This work follows a simple sequence of tests. The first obstacle is remaining competitive with, and sometimes improving over, the strong fixed grid baselines. Linear, cosine, sigmoid, and power schedules often work because they encode plausible assumptions about where a trajectory is hard to integrate. The entropy construction must do more than replace one smooth curve with another, because it must measure a path-dependent signal that fixed grids can only approximate. Cosine, sigmoid, and power grids can be good because they approximate common shapes of entropy-rate profiles. The two-dimensional sweep makes this point concrete,although it does not prove that any fixed heuristic is optimal. In those controlled experiments, cosine was often close to the entropic scheduler because its node density can approximate a boundary-heavy curve. However, this was not a separate theory of cosine schedules, and the near-tie was scenario-dependent. The result should be read as shape compatibility between cosine and the measured entropic grid, not as equivalence between the two schedules \citep{nichol2021improved,karras2022edm,lin2024common}. When the true profile is monotone, U-shaped, or boundary-heavy, these heuristics can be close. Linear grids can work when the profile is flat. The advantage of the entropy construction is that it measures the shape instead of assuming it. The same 2D runs also explain why the training-loss curves matter. In the flow-matching ideal, reducing the loss should make the learned vector field closer to the target bridge field on the training distribution. When train and validation curves decrease together, the estimated entropy-rate profile is expected to move toward the analytic U-shape. We use this only as a diagnostic for alignment between the learned field and the analytic bridge calculation, not as a convergence proof for all learned bridges. This distinction matters most at low NFE, where a few misplaced steps can determine the outcome.

%The two-dimensional experiments are the first learned test of that signal. They sit between the analytic Brownian bridge and the high-dimensional samplers. Although the transports are small, the vector fields are trained rather than closed form. They test whether a boundary-heavy rate can be recovered from a learned field. In these controlled transports, the measured rates remain boundary-heavy, the entropic grids show the largest gains at low-NFE, and ODE and SDE sweeps both show that schedule choice matters. Nonetheless, the high-dimensional analysis focuses on ODE because the probability-flow ODE is the deterministic object used in the cond--marg derivation\footnote{Grid placement error is easier to isolate without stochastic sampling noise, and AlphaFlow exposes ODE sampling rather than a matched SDE interface}.

%The next obstacle is the failure of raw entropy in EDM. This is important because EDM is outside the bridge setting. In that setting,  a raw entropy profile can be mathematically meaningful and still numerically unstable when near-singular regions dominate the grid. The \textit{log1p} regularization resolves this mismatch without changing the underlying principle: the grid should follow an information-rate profile. The ODE-error argument in \Cref{app:ode_error} explains why this calibration is linked to numerical error.
% This is the point where a weaker story would break. A raw entropy profile can be mathematically meaningful and still numerically unusable when singular or near-singular regions dominate the grid. Log tempering resolves this particular failure without changing the underlying principle: the grid should follow an information-rate profile, but the profile may need calibration before it becomes a stable numerical mesh. The ODE-error argument in \Cref{app:ode_error} explains why this is plausible rather than cosmetic.

% This is the mountain part: first we beat simple low-step issue, then EDM raw entropy becomes bad, then log1p fixes it. Need avoid sounding like a storybook. But the sequence is useful....

%The bridge-versus-diffusion distinction is the final structural test. A diffusion model often has a one-endpoint noising coordinate and a marginal score. A Schr\"odinger bridge has two endpoint constraints and a conditional path distribution. Treating the bridge as an ordinary diffusion process loses the conditional--marginal structure; treating it as unrelated to diffusion throws away useful probability-flow tools. The useful middle ground is \Cref{eq:cm_identity}. The conditional term captures endpoint-specific expansion or contraction. The marginal term captures population expansion or contraction. Their difference is the rate at which the current state resolves the bridge condition.

%This distinction also explains why a single-field estimate can be misleading. A one-field divergence can appear flat even when the bridge calculation predicts endpoint structure, because it has already collapsed the conditional and marginal roles into one object. AlphaFlow is the closer high-dimensional bridge-like test for this part of the story. It gives enough field access for the cond--marg estimator, and its estimated rate is boundary-heavy, as the Gaussian bridge calculation predicts. However, the best endpoint pLDDT is still obtained by a power grid in the available sweeps. This does not contradict the estimator result because the estimator claim concerns schedule geometry rather than endpoint pLDDT. It may instead reflect protein-specific entropy structure: protein length, flexible regions, disordered regions, and disorder-prone sequence composition can change the endpoint profile seen by the sampler \citep{bioemu2024, peptron2025}. A full case analysis by protein size, disorder content, and residue composition is beyond the scope of this work. The CAMEO trajectory pLDDT panel in \Cref{fig:cameo_alphaflow_plddt_appendix} gives a snapshot view of this protein-specific behavior by size bucket. Therefore, we treat this as a future protein-generation study, especially for AlphaFlow/Proteina-family models trained with the entropy criterion and sampled with the entropic grid \citep{jing2024alphafoldmeetsflowmatching, proteina2025}.

%The grid measures where the sampler spends integration effort, while endpoint pLDDT also depends on the training objective, model calibration, and the protein confidence metric. Once these roles are separated, the evidence has a bounded interpretation: the 2D models support the pre-flight geometry, EDM shows why mismatch requires tempering, AlphaFlow shows the conditional--marginal bridge signal, and endpoint metrics show why training-time alignment remains future work.



\section{Related work}
\label{sec:related}

\textbf{Diffusion, flow, and bridge models.} Diffusion models and score-based samplers established continuous-time generation as a numerical integration problem \citep{ho2020ddpm,kingma2021vdm,karras2022edm}. Flow matching and stochastic interpolants later simplified training by regressing vector fields along prescribed conditional paths \citep{lipman2023flow,albergo2023building,tong2024cfm,liu2023rectified}. Schr\"odinger bridges connect this view to entropy-regularized optimal transport and stochastic control \citep{schrodinger1931,schrodinger1932,leonard2014schrodinger,debortoli2021dsb,liu2023i2sb,tong2024simulation}. Recent work expands bridges to energy-based sampling, discrete spaces, structured domains, and branched trajectories \citep{liu2025asbs,tamogashev2026dataenergy,ksenofontov2025categorical,tang2026branched,wyrwal2026topological}. Our contribution is orthogonal to these training objectives: it asks how to place inference steps once a path has been learned.

\textbf{Schedulers and fast samplers.} Classical schedules and design-space analyses remain important baselines \citep{nichol2021improved,kingma2021vdm,karras2022edm,lin2024common}. DPM-Solver, DPM-Solver++, UniPC, DEIS, and STORK change the numerical update or exploit ODE structure \citep{lu2022dpm,lu2022dpmsolverpp,zhao2023unipc,zhang2022deis,tan2026stork}. Align Your Steps, optimized-step methods, dynamic-transport scheduling, and B{\'e}zierFlow search for better timesteps or scheduler parameterizations \citep{sabour2024alignsteps,xue2024optimized,tsimpos2025lipschitz,min2026bezierflow}. These approaches can be combined with an entropy-derived grid, but they do not identify the bridge-specific conditional--marginal rate.

% \textbf{Entropy and information geometry.} Entropy-based scheduling for diffusion models was introduced by \citet{stancevic2025entropic} and further connected to information dynamics by \citet{stancevic2026information}. This work extends the information view to bridge and flow settings through \Cref{eq:cm_identity}. The main difference is the object being estimated: diffusion entropic time uses a one-endpoint entropy signal, while a bridge exposes a contrast between conditional and marginal divergence. This is why the Brownian bridge calculation yields a symmetric U-shape rather than a generic diffusion-time profile.

% \textbf{Entropy and information geometry.} Entropy-based scheduling for diffusion models was introduced by \citet{stancevic2025entropic} and further connected to information dynamics by \citet{stancevic2026information}. Entropy also appears as a training-time diversity signal in protein inverse folding, where diversity-regularized DPO increases differential entropy in ProteinMPNN log-probability space \citep{park2024peptidedpo}. This work extends the information view to bridge and flow settings through \Cref{eq:cm_identity}. The main difference is the object being estimated: diffusion entropic time uses a one-endpoint entropy signal, while a bridge exposes a contrast between conditional and marginal divergence. This is why the Brownian bridge calculation yields a symmetric U-shape rather than a generic diffusion-time profile.

\textbf{Entropy and information geometry.} Entropy-based scheduling for diffusion models was introduced by \citet{stancevic2025entropic} and further connected to information dynamics by \citet{stancevic2026information}. Entropy and mutual information also appear widely as learning objectives or regularizers: entropy-regularized optimal transport and maximum-entropy control use entropy to shape couplings or policies \citep{cuturi2013sinkhorn,ziebart2008maximum,haarnoja2018soft}; information bottlenecks, InfoMax, InfoGAN, neural mutual-information estimators, and Mutual Information Machines use mutual information as a compression, representation, or generative-model objective \citep{tishby2000information,bell1995infomax,chen2016infogan,belghazi2018mine,livne2019mim}; and protein inverse folding has used diversity-regularized DPO to increase differential entropy in ProteinMPNN log-probability space \citep{park2024peptidedpo}. This work extends the information view to bridge and flow settings through \Cref{eq:cm_identity}. The main difference is the object being estimated: diffusion entropic time uses a one-endpoint entropy signal, while a bridge exposes a contrast between conditional and marginal divergence; \Cref{tab:stancevic_appendix} summarizes the comparison.

\textbf{Low-step training and scientific domains.} Distillation compresses many steps into a trained few-step model \citep{gushchin2025ibmd,he2024cdbm,ai2026jointdistillation,passaro2026stochasticfewstep}. Entropic discretization is training-free and can be used with pretrained samplers. In scientific domains, flow and bridge methods increasingly model constrained objects such as proteins, molecules, materials, and distributions over distributions \citep{jing2024alphafoldmeetsflowmatching,proteina2025,bioemu2024,irwin2024molflow,miller2024flowmm,haviv2025wasserstein}. These domains motivate measuring the path geometry directly because heuristic grids may miss endpoint stiffness or manifold constraints. A compact comparison is given in \Cref{app:comparison}.
\section{Limitations and future work}
\label{sec:limitations}

%This work does not train a high-dimensional model with the entropy criterion. This is a real limitation and a useful boundary around the claims. The present method is designed for inference-time use under low compute budgets. Although the two-dimensional models are trained vector-field probes, they are controlled pre-flight checks rather than a replacement for larger-scale entropy-objective training. Therefore, a natural next step is to train bridge or flow models with an explicit cond--marg entropy criterion and then use the corresponding entropic scheduler at inference time.

Overall, the results support conditional--marginal entropy rate as a practical low-NFE allocation signal, provided it is regularized into a stable grid and interpreted separately from endpoint-quality metrics. The estimator depends on access to conditional and marginal fields. Some pretrained models expose only a single field or require a surrogate for one of the terms. In such cases the resulting grid should be treated as a diagnostic unless both terms are matched to the theory. Hutchinson estimation introduces variance, and raw boundary singularities can be too sharp for numerical stability. Log tempering is a practical calibration, not a theorem. The closed-form U-shape is also a Gaussian Brownian-bridge result. For non-Gaussian learned bridges, \Cref{eq:cm_identity} remains the target, but the shape of the rate is an empirical or model-specific question. While the current results are encouraging, optimized entropic scheduling for protein generation deserves a separate research program. Short proteins, long proteins, flexible loops, and disordered regions may have different entropy profiles. Models such as AlphaFlow, BioEmu, Proteina, and Peptron make it possible to study those profiles as biological signals rather than as nuisance quantities \citep{jing2024alphafoldmeetsflowmatching,bioemu2024,proteina2025, peptron2025}. Besides protein generation, the same question could be tested in image generation by training or calibrating EDM- or latent-diffusion models with the entropy criterion, and then sampling them with an entropic grid at inference time \citep{karras2022edm,rombach2022ldm,ho2020ddpm}. In both domains, the next question is whether training and inference can be aligned around the same cond--marg information rate.

%\section{Conclusion}
%\label{sec:conclusion}

%Three strands converge in this paper: the continuity-equation identity that turns entropy rate into an expected divergence, the Brownian-bridge calculation that turns this identity into a concrete U-shaped conditional profile, and the high-dimensional evidence showing how the same idea behaves in AlphaFlow and EDM. Taken together, these strands give a principled but bounded conclusion. Entropy is a guide for low-budget generative sampling when it is estimated as the bridge information rate, calibrated as a numerical grid, and interpreted separately from endpoint confidence.


% \begin{ack}
% Use unnumbered first level headings for the acknowledgments. All acknowledgments
% go at the end of the paper before the list of references. Moreover, you are required to declare
% funding (financial activities supporting the submitted work) and competing interests (related financial activities outside the submitted work).
% More information about this disclosure can be found at: \url{https://neurips.cc/Conferences/2026/PaperInformation/FundingDisclosure}.
%
% Do {\bf not} include this section in the anonymized submission, only in the final paper. You can use the \texttt{ack} environment provided in the style file to automatically hide this section in the anonymized submission.
% \end{ack}


 \bibliographystyle{plainnat}
  % \bibliographystyle{unsrtnat}
\bibliography{references}


%%%%%%%%%%%%%%%%%%%%%%%%%%%%%%%%%%%%%%%%%%%%%%%%%%%%%%%%%%%%

\appendix

\newpage
\begin{center}
{\Large\bfseries Appendix}
\end{center}
\vspace{1em}

\section{A guided derivation of the conditional--marginal rate}
\label{app:full_derivations}

\paragraph{Marginal entropy rate.}
\label{app:marginal_entropy}

We assumed in the main text that the density and vector field are smooth enough for integration by parts and that boundary terms vanish. We now walk through the calculation. Consider first the marginal entropy
\begin{equation}
  \Ent(X_t)=-\int p_t(\vx)\log p_t(\vx)\,\dd \vx,
\end{equation}
and differentiate with respect to time:
\begin{align}
  \frac{\dd}{\dd t}\Ent(X_t)
  &=
  -\int \partial_t p_t(\vx)\log p_t(\vx)\,\dd\vx
  -\int p_t(\vx)\frac{\partial_t p_t(\vx)}{p_t(\vx)}\,\dd\vx \\
  &=
  -\int \partial_t p_t(\vx)\log p_t(\vx)\,\dd\vx
  -\frac{\dd}{\dd t}\int p_t(\vx)\,\dd\vx \\
  &=
  -\int \partial_t p_t(\vx)\log p_t(\vx)\,\dd\vx.
\end{align}
The second term vanishes because $p_t$ remains normalized. Next, use the continuity equation,
\begin{align}
  \frac{\dd}{\dd t}\Ent(X_t)
  &=
  \int \diver(p_t(\vx)\bar{\vv}_t(\vx))\log p_t(\vx)\,\dd\vx \\
  &=
  -\int p_t(\vx)\bar{\vv}_t(\vx)^\top \nabla_{\vx}\log p_t(\vx)\,\dd\vx \\
  &=
  -\int \bar{\vv}_t(\vx)^\top \nabla_{\vx}p_t(\vx)\,\dd\vx \\
  &=
  \int p_t(\vx)\diver \bar{\vv}_t(\vx)\,\dd\vx \\
  &=
  \E_{X_t}\!\left[\diver\bar{\vv}_t(X_t)\right].
\end{align}

\paragraph{Conditional entropy rate.}
\label{app:conditional_entropy}

The conditional calculation repeats the same story while holding the bridge condition fixed. For each condition $Z=z$, the conditional density follows its own continuity equation, so
\begin{equation}
  \frac{\dd}{\dd t}\Ent(X_t\mid Z=z)
  =
  \E_{X_t\mid Z=z}\!\left[\diver \vv_t(X_t\mid z)\right].
\end{equation}
Averaging over $Z$ gives
\begin{equation}
  \frac{\dd}{\dd t}\Ent(X_t\mid Z)
  =
  \E_{Z,X_t\mid Z}\!\left[\diver \vv_t(X_t\mid Z)\right].
\end{equation}
Now take into consideration the entropy chain rule,
\begin{equation}
  \Ent(Z\mid X_t)=\Ent(X_t\mid Z)+\Ent(Z)-\Ent(X_t),
\end{equation}
and subtract the marginal entropy rate derived above. We obtain
\begin{equation}
  \frac{\dd}{\dd t}\Ent(Z\mid X_t)
  =
  \E_{Z,X_t\mid Z}\!\left[\diver \vv_t(X_t\mid Z)\right]
  -
  \E_{X_t}\!\left[\diver \bar{\vv}_t(X_t)\right].
\end{equation}

% Draft scratch:
% Alternative derivation path maybe: start from H[Z|X_t] directly and differentiate p(z|x_t). This becomes messier because posterior changes with time. Chain rule is cleaner and probably what we should keep.
% \begin{align}
%   \frac{\dd}{\dd t}H(Z|X_t)
%   &= -\frac{\dd}{\dd t}\int p_t(x,z)\log p_t(z|x)\dd x\dd z.
% \end{align}
% This is possible but not the friendliest route.

\paragraph{Score representation.}
\label{app:score_derivation}

The divergence form is the estimator used in this paper, but the score form clarifies what is being measured. The marginal field is
\begin{equation}
  \bar{\vv}_t(\vx)=\int p_t(z\mid \vx)\vv_t(\vx\mid z)\,\dd z.
\end{equation}
Taking divergence,
\begin{align}
  \diver \bar{\vv}_t(\vx)
  &=
  \int p_t(z\mid\vx)\diver \vv_t(\vx\mid z)\,\dd z
  +
  \int \nabla_{\vx}p_t(z\mid\vx)^\top \vv_t(\vx\mid z)\,\dd z.
\end{align}
Multiply by $p_t(\vx)$ and integrate over $\vx$. The first term becomes the expected conditional divergence. For the second term, Bayes' rule gives
\begin{equation}
  \nabla_{\vx}\log p_t(z\mid\vx)
  =
  \nabla_{\vx}\log p_t(\vx\mid z)
  -
  \nabla_{\vx}\log p_t(\vx).
\end{equation}
Therefore
\begin{align}
  \E_{X_t}\!\left[\diver\bar{\vv}_t(X_t)\right]
  &=
  \E_{Z,X_t\mid Z}\!\left[\diver\vv_t(X_t\mid Z)\right] \nonumber \\
  &\quad+
  \E_{Z,X_t\mid Z}\!\left[
  (\nabla\log p_t(X_t\mid Z)-\nabla\log p_t(X_t))^\top
  \vv_t(X_t\mid Z)
  \right].
\end{align}
Rearranging yields \Cref{eq:score_form}. This shows that the rate measures the conditional field against the gap between the conditional score and the marginal score.

\paragraph{Gaussian conditional vector field.}
\label{app:gaussian_derivation}

The Brownian-bridge calculation in the main text is a special case of a Gaussian conditional path. Let
\begin{equation}
  X_t=\mu(z,t)+\sigma(z,t)\veps,\qquad \veps\sim\mathcal{N}(0,I).
\end{equation}
For fixed noise $\veps$,
\begin{equation}
  \partial_t X_t=\partial_t\mu(z,t)+\partial_t\sigma(z,t)\veps.
\end{equation}
Since $\veps=(X_t-\mu(z,t))/\sigma(z,t)$,
\begin{equation}
  \vv_t(X_t\mid z)=\partial_t\mu(z,t)+\partial_t\sigma(z,t)\frac{X_t-\mu(z,t)}{\sigma(z,t)}.
\end{equation}
The Gaussian score is
\begin{equation}
  \nabla_{X_t}\log p_t(X_t\mid z)=-\frac{X_t-\mu(z,t)}{\sigma^2(z,t)}.
\end{equation}
Substitution gives
\begin{equation}
  \vv_t(X_t\mid z)
  =
  \partial_t\mu(z,t)
  -
  \sigma(z,t)\partial_t\sigma(z,t)
  \nabla_{X_t}\log p_t(X_t\mid z).
\end{equation}

\section{Brownian bridge details}
\label{app:bb_details}

\subsection{Probability-flow field}

For $m_t=(1-t)x_0+tx_1$ and $\sigma(t)=\sigma_0\sqrt{t(1-t)}$,
\begin{equation}
  \partial_t m_t=x_1-x_0.
\end{equation}
The time derivative of the standard deviation is
\begin{align}
  \partial_t\sigma(t)
  &=
  \sigma_0\frac{1}{2\sqrt{t(1-t)}}(1-2t),
\end{align}
so
\begin{equation}
  \sigma(t)\partial_t\sigma(t)
  =
  \sigma_0\sqrt{t(1-t)}
  \frac{\sigma_0(1-2t)}{2\sqrt{t(1-t)}}
  =
  \frac{\sigma_0^2(1-2t)}{2}.
\end{equation}
The conditional score is
\begin{equation}
  \nabla_{\vx}\log p_t(\vx\mid x_0,x_1)
  =
  -\frac{\vx-m_t}{\sigma_0^2t(1-t)}.
\end{equation}
Substituting into the Gaussian field,
\begin{align}
  \vv_t(\vx\mid x_0,x_1)
  &=
  x_1-x_0
  -
  \frac{\sigma_0^2(1-2t)}{2}
  \left[-\frac{\vx-m_t}{\sigma_0^2t(1-t)}\right] \\
  &=
  x_1-x_0
  +
  \frac{1-2t}{2t(1-t)}(\vx-m_t).
\end{align}

\subsection{SDE drift and the factor of two}
\label{app:sde_ode_details}

The Brownian-bridge SDE drift is
\begin{equation}
  u^o_t(\vx\mid x_0,x_1)
  =
  x_1-x_0
  +
  \frac{1-2t}{t(1-t)}(\vx-m_t).
\end{equation}
Comparing with the probability-flow field gives
\begin{equation}
  u^o_t(\vx\mid x_0,x_1)
  =
  2\vv_t(\vx\mid x_0,x_1)-(x_1-x_0).
\end{equation}
The factor of two arises from the chain-rule term $\sigma(t)\partial_t\sigma(t)$ in the probability-flow derivation. If one applies the marginal probability-flow transformation to the conditional objects, the score correction cancels the bridge contraction term incorrectly:
\begin{align}
  u^o_t+\sigma_0^2(1-2t)\nabla_{\vx}\log p_t(\vx\mid x_0,x_1)
  &=
  x_1-x_0
  +
  \frac{1-2t}{t(1-t)}(\vx-m_t)
  -
  \frac{1-2t}{t(1-t)}(\vx-m_t) \\
  &=
  x_1-x_0.
\end{align}
The result differs from \Cref{eq:bb_ode_field}. The probability-flow relation is valid for marginal fields after endpoint mixing, not for this conditional conversion.

\section{Interpolation profiles and bridge specificity}
\label{app:interpolation_profiles}

We raised an important scope question: is the U-shape a generic property of every interpolation, or a bridge-specific calculation? The answer is bridge-specific. Different interpolants produce different analytic divergence profiles, and the learned field can deviate from the exact conditional formula. This is why the main text treats the Brownian-bridge profile as an anchor and the conditional--marginal estimator as the general object.

\begin{table}[H]
  \caption{Analytic divergence profiles for common conditional interpolants. Here $d$ is the ambient dimension. The table illustrates that different path choices lead to different schedule signals.}
  \label{tab:interpolation_profiles}
  \centering
  \small
  \begin{tabular}{P{0.24\linewidth}P{0.42\linewidth}P{0.22\linewidth}}
    \toprule
    Interpolation & Exact conditional divergence & Qualitative profile \\
    \midrule
    Linear optimal-transport interpolation & $0$ for the exact linear conditional field & Flat before learning error \\
    Variance-preserving diffusion path & $d\,t/(1-t^2)$ under the standard VP parameterization & Monotone boundary growth \\
    Cosine diffusion path & $d\pi\tan(\pi t/2)/2$ under the cosine noise parameterization & Monotone boundary growth \\
    Brownian bridge & $d(1-2t)/[2t(1-t)]$ for the probability-flow field & Symmetric U-shape in magnitude \\
    \bottomrule
  \end{tabular}
\end{table}

For Schr\"odinger bridges represented as mixtures of Brownian bridges, the conditional Brownian divergence in \Cref{eq:bb_divergence} does not depend on the endpoint pair $(x_0,x_1)$ or on the spatial position $\vx$. Therefore
\begin{equation}
  \E_{(X_0,X_1),X_t\mid X_0,X_1}
  \left[\diver\vv_t(X_t\mid X_0,X_1)\right]
  =
  d\,\frac{1-2t}{2t(1-t)}
\end{equation}
for the conditional term of the Brownian reference bridge. The marginal term in \Cref{eq:cm_identity} remains model- and coupling-dependent, which is why the full bridge rate is not reduced to the conditional term alone.

\section{Boundary concentration ratio}
\label{app:bcr}

We introduced boundary concentration ratio after observing that many useful schedules differ mainly in where they place density near $t=0$ and $t=1$. The entropy rate gives a full curve, but comparisons across many schedules are easier when one also has a scalar diagnostic. The diagnostic should answer a narrow question: how much more schedule mass lies near the two endpoints than would lie there under a uniform grid?

Take into consideration a normalized schedule density $q(t)$ on $[0,1]$ with $\int_0^1q(t)\dd t=1$. For a boundary width $\epsilon\in(0,1/2)$, define the endpoint mass
\begin{equation}
  M_\epsilon(q)
  =
  \int_0^\epsilon q(t)\,\dd t+\int_{1-\epsilon}^1 q(t)\,\dd t.
\end{equation}
The same two endpoint intervals have mass $2\epsilon$ under a uniform grid density. We therefore define
\begin{equation}
  \bcr_\epsilon(q)
  =
  \frac{M_\epsilon(q)}{2\epsilon}.
  \label{eq:bcr}
\end{equation}
A uniform grid has $\bcr_\epsilon=1$. Values larger than one mean that the schedule assigns more node density to the boundary windows than uniform spacing. Values below one mean that the schedule avoids the endpoints relative to the uniform grid.

For a discrete grid, one can estimate $q$ from the interval widths or from the inverse-CDF density used to construct the grid. The density-based version is preferable because it is less sensitive to whether the endpoints themselves are counted as nodes. The assumptions behind BCR are simple: the same $\epsilon$ must be used for all schedules, the grid must be monotone, and the diagnostic should be computed on the schedule density rather than on downstream metrics.

BCR is useful beyond this paper because it separates a geometric property of the grid from sample quality. A high BCR can indicate endpoint stiffness, bridge contraction, or a raw entropy singularity. It can also warn that tempering may be needed, since over-resolving a singular endpoint can waste low-NFE budget. BCR is therefore a schedule diagnostic and not a main claim. A model can have high boundary concentration and still perform poorly if the estimator is mismatched or if the training objective favors another grid.

\section{ODE error derivation}
\label{app:ode_error}

For an order-$p$ one-step method, suppose the local truncation error on a step of width $h$ is bounded by
\begin{equation}
  \|e_{k+1}\|\le C(t_k)h_k^{p+1}+O(h_k^{p+2}).
\end{equation}

% Draft scratch:
% This is only a proxy. We should not say entropy is the actual LTE constant. The honest thing is: divergence can point to stiffness / expansion / compression and that can correlate with where the solver needs smaller steps.
% Extra connection maybe:
% \begin{equation}
%   |\mathrm{tr}(J)| \le d\|J\|_2,
%   \qquad
%   \|J\|_2 \ge |\mathrm{tr}(J)|/d.
% \end{equation}
% So large divergence implies at least one large Jacobian direction, but small divergence does not mean non-stiff because eigenvalues can cancel.
Let $\rho(t)$ be a continuous density of grid points. Locally $h(t)\approx 1/(N\rho(t))$. Ignoring constants independent of $\rho$, a continuous proxy for total error is
\begin{equation}
  \mathcal{E}[\rho]=\int_0^1 C(t)\rho(t)^{-p}\,\dd t,
  \qquad
  \int_0^1\rho(t)\,\dd t=1.
\end{equation}
The Lagrangian is
\begin{equation}
  \mathcal{L}[\rho]
  =
  \int_0^1 C(t)\rho(t)^{-p}\,\dd t
  +
  \lambda\left(\int_0^1\rho(t)\,\dd t-1\right).
\end{equation}
The stationarity condition is
\begin{equation}
  -pC(t)\rho(t)^{-p-1}+\lambda=0,
\end{equation}
which gives
\begin{equation}
  \rho(t)\propto C(t)^{1/(p+1)}.
\end{equation}
Entropy-rate scheduling uses $|\dot{\Ent}(Z\mid X_t)|$ as a measurable proxy for $C(t)$.

\section{Algorithms}
\label{app:algorithms}

This appendix states the inference-time pipeline as algorithms. The algorithms do not assume a particular implementation. They separate the calibration stage, where the entropy-rate curve is estimated once, from the sampling stage, where all schedulers are compared at matched NFE.

\begin{algorithm}[H]
  \caption{Calibration Set and Time Mesh Construction}
  \label{alg:calibration}
  \begin{algorithmic}[1]
    \REQUIRE Calibration conditions $\{z_i\}_{i=1}^n$, time range $[\epsilon,1-\epsilon]$, number of mesh points $M$
    \ENSURE Calibration states $\{x_{ij}\}$ and paired conditions $\{z_i,s_j\}$
    \STATE Choose a monotone mesh $0<\epsilon=s_1<\cdots<s_M=1-\epsilon<1$
    \FOR{$i=1,\ldots,n$}
      \STATE Draw or retrieve the bridge condition $z_i$
      \FOR{$j=1,\ldots,M$}
        \STATE Construct the state $x_{ij}\sim p_{s_j}(\cdot\mid z_i)$ when the conditional simulator is available
        \STATE Otherwise, use a model trajectory state paired with its condition and time
      \ENDFOR
    \ENDFOR
    \RETURN Paired calibration set $\{(x_{ij},z_i,s_j)\}$
  \end{algorithmic}
\end{algorithm}

\begin{algorithm}[H]
  \caption{Hutchinson Divergence of a Vector Field}
  \label{alg:hutchinson}
  \begin{algorithmic}[1]
    \REQUIRE Vector field $f_t$, state $x$, time $t$, probes $\{\vu_\ell\}_{\ell=1}^m$ with $\E[\vu_\ell\vu_\ell^\top]=I$
    \ENSURE Estimate of $\diver f_t(x)$
    \STATE $a \leftarrow 0$
    \FOR{$\ell=1,\ldots,m$}
      \STATE Evaluate the Jacobian-vector product $w_\ell \leftarrow \nabla_x f_t(x)\vu_\ell$
      \STATE $a \leftarrow a + \vu_\ell^\top w_\ell$
    \ENDFOR
    \RETURN $a/m$
  \end{algorithmic}
\end{algorithm}

\begin{algorithm}[H]
  \caption{Conditional--Marginal Entropy-Rate Estimation}
  \label{alg:cm_estimator}
  \begin{algorithmic}[1]
    \REQUIRE Calibration set $\{(x_{ij},z_i,s_j)\}$, conditional field $\vv_t(\cdot\mid z)$, marginal field $\bar{\vv}_t$, probes per state $m$
    \ENSURE Estimated rate values $\{\hat r(s_j)\}_{j=1}^M$ and uncertainty estimates
    \FOR{$j=1,\ldots,M$}
      \STATE $d_{j1},\ldots,d_{jn} \leftarrow$ empty
      \FOR{$i=1,\ldots,n$}
        \STATE Draw shared probes $\{\vu_{\ell ij}\}_{\ell=1}^m$
        \IF{analytic conditional divergence is available}
          \STATE $c_{ij}\leftarrow \diver\vv_{s_j}(x_{ij}\mid z_i)$ from the analytic formula
        \ELSE
          \STATE $c_{ij}\leftarrow$ \Cref{alg:hutchinson} applied to $\vv_{s_j}(\cdot\mid z_i)$ at $x_{ij}$
        \ENDIF
        \STATE $m_{ij}\leftarrow$ \Cref{alg:hutchinson} applied to $\bar{\vv}_{s_j}$ at $x_{ij}$ with the same probes
        \STATE $d_{ji}\leftarrow c_{ij}-m_{ij}$
      \ENDFOR
      \STATE $\hat r(s_j)\leftarrow \left|n^{-1}\sum_{i=1}^n d_{ji}\right|$
      \STATE $\widehat{\mathrm{se}}(s_j)\leftarrow$ standard error of $\{d_{ji}\}_{i=1}^n$
    \ENDFOR
    \RETURN Rate curve $\hat r$ and uncertainty curve $\widehat{\mathrm{se}}$
  \end{algorithmic}
\end{algorithm}

\begin{algorithm}[H]
  \caption{Rate Postprocessing and Grid Construction}
  \label{alg:grid}
  \begin{algorithmic}[1]
    \REQUIRE Mesh $\{s_j\}_{j=1}^M$, estimated rate $\hat r$, transform $\phi\in\{r,\log(1+r)\}$, sampling steps $N$
    \ENSURE Monotone grid $0=t_0<t_1<\cdots<t_N=1$
    \STATE Replace non-finite values by local interpolation and clip $\hat r$ to a small positive floor
    \STATE Smooth the clipped curve only on the calibration mesh, preserving endpoint ordering
    \STATE $a_j\leftarrow \phi(\hat r(s_j))$ for $j=1,\ldots,M$
    \STATE Normalize $a_j$ into a density $q_j$ by numerical quadrature
    \STATE Compute the cumulative curve $Q(s_j)\approx\int_0^{s_j}q(u)\dd u$
    \STATE Enforce monotonicity of $Q$ and set $Q(0)=0$, $Q(1)=1$
    \FOR{$k=0,\ldots,N$}
      \STATE Set $t_k\leftarrow Q^{-1}(k/N)$ by monotone interpolation
    \ENDFOR
    \RETURN Grid $\{t_k\}_{k=0}^N$ and density summary $q$
  \end{algorithmic}
\end{algorithm}

\begin{algorithm}[H]
  \caption{Matched-NFE Evaluation Protocol}
  \label{alg:evaluation}
  \begin{algorithmic}[1]
    \REQUIRE Model sampler, scheduler family $\mathcal{S}$, NFE budgets $\mathcal{N}$, seeds $\mathcal{B}$, evaluation metrics $\mathcal{M}$
    \ENSURE Matched comparison table across schedulers
    \FOR{scheduler $S\in\mathcal{S}$}
      \FOR{budget $N\in\mathcal{N}$}
        \STATE Build the grid for $S$ with exactly $N$ intervals
        \FOR{seed $b\in\mathcal{B}$}
          \STATE Generate samples with the same model, solver family, and NFE budget
          \STATE Evaluate all metrics in $\mathcal{M}$ on the generated samples
        \ENDFOR
      \ENDFOR
    \ENDFOR
    \STATE Aggregate means, standard deviations, and paired deltas relative to selected baselines
    \RETURN Matched-NFE metric table and per-sample records
  \end{algorithmic}
\end{algorithm}

\begin{algorithm}[H]
  \caption{Paired Bootstrap Confidence Intervals}
  \label{alg:bootstrap}
  \begin{algorithmic}[1]
    \REQUIRE Per-unit metric records for scheduler $A$ and baseline $B$, bootstrap draws $R$, confidence level $1-\alpha$
    \ENSURE Mean paired difference and percentile interval
    \STATE Match records by evaluation unit, seed, dataset split, and NFE budget
    \FOR{$r=1,\ldots,R$}
      \STATE Resample matched units with replacement
      \STATE Compute the paired difference $\Delta_r=\mathrm{metric}(A)-\mathrm{metric}(B)$ on the resample
    \ENDFOR
    \STATE Report $\bar{\Delta}$ on the original matched records
    \STATE Report percentiles $\alpha/2$ and $1-\alpha/2$ of $\{\Delta_r\}_{r=1}^R$
    \RETURN Paired estimate and confidence interval
  \end{algorithmic}
\end{algorithm}

\begin{algorithm}[H]
  \caption{Boundary-Concentration Diagnostic}
  \label{alg:bcr}
  \begin{algorithmic}[1]
    \REQUIRE Schedule density $q(t)$ or grid density estimate, boundary width $\epsilon\in(0,1/2)$
    \ENSURE Boundary concentration ratio $\bcr_\epsilon$
    \STATE Estimate boundary mass $M_\epsilon=\int_0^\epsilon q(t)\dd t+\int_{1-\epsilon}^1q(t)\dd t$
    \STATE Normalize by the uniform-grid boundary mass $2\epsilon$
    \STATE $\bcr_\epsilon\leftarrow M_\epsilon/(2\epsilon)$
    \RETURN $\bcr_\epsilon$
  \end{algorithmic}
\end{algorithm}

\section{Computational notes}
\label{app:compute}

\begin{table}[H]
  \caption{Complexity of grid construction. Sampling cost at fixed NFE is unchanged after the grid has been constructed.}
  \label{tab:complexity}
  \centering
  \small
  \begin{tabular}{P{0.22\linewidth}P{0.24\linewidth}P{0.22\linewidth}P{0.20\linewidth}}
    \toprule
    Method & Calibration compute & Extra memory & TFLOPS/byte note \\
    \midrule
    Analytic bridge rate & $O(M)$ scalar evaluations & $O(N)$ grid & Negligible \\
    Exact divergence & $O(Mnd)$ derivative products & Large Jacobian or repeated products & Usually impractical \\
    Hutchinson rate & $O(Mnm)$ derivative-vector products & Probe buffers plus activations & Kernel-dependent \\
    Reused entropy grid & None during sampling & $O(N)$ grid & Same NFE as baseline \\
    \bottomrule
  \end{tabular}
\end{table}

Here $M$ is the number of calibration times, $n$ the number of calibration states, $m$ the number of Hutchinson probes, $d$ the dimension, and $N$ the number of sampling steps. Hardware-level TFLOPS/byte depends on kernels, precision, and hardware, so the robust comparison is symbolic.

\section{Hypotheses and conclusions}
\label{app:hypotheses}

We guided the work with three overall questions. The first asks whether the bridge calculation predicts a distinctive information geometry. The second asks whether learned two-dimensional probes preserve that geometry after the vector field is trained. The third asks whether cond--marg estimator and its grid help interpret high-dimensional low-NFE behavior without turning an allocation diagnostic into an endpoint-quality claim.

% We guided the work with three overall questions. The first asks whether the bridge calculation predicts a distinctive information geometry. The second asks whether a conditional--marginal estimator can reveal that geometry in high-dimensional samplers. The third asks whether the resulting grid improves or explains low-NFE behavior without turning an allocation diagnostic into an endpoint-quality claim.

\begin{table}[H]
  \caption{Questions that guided the evidence and the corresponding conclusions.}
  \label{tab:hypotheses}
  \centering
  \small
  \begin{tabular}{P{0.28\linewidth}P{0.42\linewidth}P{0.18\linewidth}}
    \toprule
    Guiding question & Evidence & Conclusion \\
    \midrule
    What quantity should a bridge schedule measure? & \Cref{eq:cm_identity} and \Cref{eq:bb_divergence} show that the bridge rate is a conditional--marginal divergence contrast, with a U-shaped Brownian-bridge conditional term. & Use a two-term bridge rate, not a single-field proxy. \\
    Does this geometry appear in high dimension? & AlphaFlow conditional--marginal profiles show a boundary-heavy U-shape, while EDM gives a different profile and therefore a useful mismatch case. & The rate is model-dependent and must be measured or derived. \\
    Can the measured rate guide low-NFE sampling? & EDM five-step FID improves with log-tempered entropy, and AlphaFlow allocation diagnostics align with the bridge calculation while endpoint pLDDT remains training-objective dependent. & Entropy is a grid signal, not a universal endpoint-quality guarantee. \\
    \bottomrule
  \end{tabular}
\end{table}


\section{Additional method comparison}
\label{app:comparison}

\begin{table}[H]
  \caption{Comparison with related acceleration and scheduling families.}
  \label{tab:comparison}
  \centering
  \small
  \begin{tabular}{P{0.24\linewidth}P{0.30\linewidth}P{0.32\linewidth}}
    \toprule
    Family & Main object & Relation to this work \\
    \midrule
    Heuristic grids & Fixed analytic time maps & Strong baselines, but not path-measured \\
    Align Your Steps & Optimized discrete schedule & Empirical optimization rather than entropy identity \\
    Lipschitz and transport scheduling & Smoothness or transport-cost proxy & Related to error constants, no conditional--marginal bridge decomposition \\
    Fast ODE solvers & Numerical update formula & Complementary to grid choice \\
    Distillation & New few-step model & Requires training, unlike inference-time grid construction \\
    Bridge and flow coupling & Training path geometry & Changes the learned path rather than measuring the grid density \\
    \bottomrule
  \end{tabular}
\end{table}

    \begin{table}[H]
      \caption{Closest entropy-based scheduler comparison. The bridge-specific change is the conditional--marginal divergence contrast.}
      \label{tab:stancevic_appendix}
      \centering
      \small
      \setlength{\tabcolsep}{3pt}
      \begin{tabular}{@{}P{0.14\linewidth}@{\hspace{0.02\linewidth}}P{0.39\linewidth}@{\hspace{0.02\linewidth}}P{0.39\linewidth}@{}}
        \toprule
        Aspect & Entropic diffusion time & Conditional--marginal bridge time \\
        \midrule
        Object & One-endpoint diffusion process & Conditional bridge or flow path \\
        Signal & Entropy of data given noised state & $\dot{\Ent}(Z\mid X_t)$ from conditional minus marginal divergence \\
        Estimator & Diffusion loss or single-field formula & Matched conditional and marginal divergence estimates \\
        Shape & Often asymmetric in noise time & Symmetric U-shape for Brownian bridges \\
        Use & Inference-time reparameterization & Inference-time grid for bridge and flow samplers \\
        \bottomrule
      \end{tabular}
    \end{table}

\section{Protein endpoint confidence metrics}
\label{app:plddt_metric}

Protein endpoint quality is reported with predicted local distance difference test (plDDT), a per-residue confidence estimate associated with distance agreements \citep{jumper2021, mariani2013lddt}. The score is usually reported on a 0--100 scale. Higher values indicate higher local confidence. We average plDDT over generated endpoint structures and use it as an endpoint confidence \emph{proxy}. This is useful because it is cheap, standardized, and sensitive to local structural plausibility. However, it is not the same object as the cond--marg entropy rate. plDDT does not measure where the sampler placed time steps, and it can depend on the target protein, sequence length, flexible or disordered regions, model calibration, and the evaluator. Hence, pLDDT is interpreted together with diversity and trajectory diagnostics, not as direct proof that one grid is the correct entropy grid.

\paragraph{Protein evaluation funnel.}
\label{app:protein_eval_funnel}
CAMEO22 and ATLAS are used here as evaluation sources, not as evidence that the entropy grid was optimized during AlphaFlow training. The Sankey diagram in \Cref{fig:alphaflow_testing_sankey_appendix} separates the available benchmark proteins from the subset that passed the native-score testing path. The count table behind the figure reports 48 available ATLAS proteins and 72 available CAMEO proteins; 115 proteins were tested after filtering (45 ATLAS, 70 CAMEO), with give large proteins not tested due to system timeouts. These counts describe evaluation coverage. They should be read separately from endpoint pLDDT, which remains a confidence proxy, and separately from the cond--marg rate, which is the schedule-geomtry signal.


\section{Toy and synthetic data}
\label{app:figures}

The synthetic appendix is included to show what the theory predicts before the high-dimensional stress tests. These plots are not used as the main evidence. They help the reader see why a bridge can ask for a non-uniform grid even when the downstream model is simple.


    \begin{figure}[!t]
      \centering
      \begin{minipage}[t]{0.49\linewidth}
        \centering
        \adjustbox{max width=\linewidth, max height=0.32\textheight, valign=t}{%
          \includegraphics{entropy_tau_grid.pdf}%
        }
      \end{minipage}
      \hfill
      \begin{minipage}[t]{0.49\linewidth}
        \centering
        \adjustbox{max width=\linewidth, max height=0.32\textheight, valign=t}{%
          \includegraphics{fig_transport_scenarios_grid_tarchic_vector_field_rose_points.pdf}%
        }
      \end{minipage}
      \caption{Controlled transport setup and entropy-induced grids. Left: entropy-rate profiles and induced inverse-CDF time grids show the core mechanism: measure where the path carries information, normalize that rate, and place time nodes by the corresponding cumulative distribution. Right: synthetic transport scenarios used as controlled probes, contrasting continuous--continuous, continuous--discrete, discrete--continuous, and discrete--discrete endpoint structure.}
      \label{fig:controlled_transport_setup_appendix}
    \end{figure}

% \begin{figure}[H]
%   \centering
%   \includegraphics[width=0.78\linewidth]{fig_transport_scenarios_grid.pdf}
%   \caption{Synthetic transport scenarios used as controlled probes. The grid contrasts continuous--continuous, continuous--discrete, discrete--continuous, and discrete--discrete endpoint structure. The point of this panel is to show that the bridge path can face different endpoint constraints even before any high-dimensional neural sampler is introduced.}
%   \label{fig:transport_scenarios_appendix}
% \end{figure}

% \begin{figure}[H]
%   \centering
%   \includegraphics[width=0.78\linewidth]{entropic_condmarg_variants_comparison.pdf}
%   \caption{Synthetic entropy-rate profiles. The reader should observe that entropy production is not generally flat in time. Boundary-heavy and U-shaped profiles appear naturally in bridge-like transports, which motivates measuring a rate before committing to a grid.}
%   \label{fig:synthetic_entropy_rate}
% \end{figure}

% \begin{figure}[H]
%   \centering
%   \includegraphics[width=0.78\linewidth]{entropy_tau_grid.pdf}
%   \caption{Synthetic entropic grids obtained from the rate profiles. The figure shows how inverse-CDF allocation converts an information-rate curve into a sequence of time nodes. Dense regions in the grid correspond to high estimated entropy-rate mass.}
%   \label{fig:synthetic_tau_grid}
% \end{figure}

\begin{figure}[H]
  \centering
  \includegraphics[width=0.78\linewidth]{fig_mmd_vs_steps_ode_sde.pdf}
  \caption{Scenario-level comparison on the toy transports. The figure is included to make the synthetic evidence explicit: the schedule effect depends on the endpoint geometry, and the role of the toy data is to isolate this dependence before moving to EDM and AlphaFlow.}
  \label{fig:scenario_comparison_appendix}
\end{figure}

\begin{figure}[H]
  \centering
  \includegraphics[width=0.78\linewidth]{fig_mmd_vs_steps_ode_sde.pdf}
  \caption{Initial ODE/SDE probing on synthetic data. The main observation is that low-NFE behavior is more sensitive to grid placement than high-NFE behavior, and deterministic integration exposes this sensitivity more directly than stochastic integration.}
  \label{fig:mmd_steps_appendix}
\end{figure}

\begin{figure}[H]
  \centering
  \includegraphics[width=0.78\linewidth]{cond_marg_timeseries_10step_all_schedulers_edm_alphaflow.pdf}
  \caption{Reference scheduler shapes at $\sigma=0.5$. This plot is kept as a reference for how common heuristic grids distribute nodes relative to the bridge-inspired profiles. It is not used to claim that one heuristic is universally best.}
  \label{fig:scheduler_shapes_appendix}
\end{figure}

  \section{Evaluation Grids}
  \label{app:evaluation_grids}

% The following figures are visual checks for the grids and trajectories discussed in the main text. They are separated from endpoint metric panels to avoid conflating schedule geometry with sample-quality claims.
The following figures are grid and trajectory diagnostics referenced from the main text.

\begin{figure}[t]
  \centering
  \begin{subfigure}[t]{0.49\linewidth}
    \centering
    \includegraphics[width=\linewidth]{figures__EDM_trajectory_grid_ode_heun_10steps_linear.pdf}
    \caption{Linear.}
  \end{subfigure}
  \hfill
  \begin{subfigure}[t]{0.49\linewidth}
    \centering
    \includegraphics[width=\linewidth]{figures__EDM_trajectory_grid_ode_heun_10steps_entropic_log1p.pdf}
    \caption{Cond--marg log1p.}
  \end{subfigure}
  \caption{EDM ten-step ODE-Heun trajectory grids. The quantitative five-step claim is carried by \Cref{tab:edm-fid-ode}; Log1p allocation is visually close to linear while remaining a distinct measured grid.}
  \label{fig:edm_images}
\end{figure}


% \begin{figure}[H]
%   \centering
%   \includegraphics[width=0.82\linewidth]{EDM_trajectory_grid_ode_heun_10steps.pdf}
%   \caption{EDM ODE-Heun trajectory grid at 10 steps. Compared with the five-step main-text figure, this panel shows how schedule differences become less visually pronounced as NFE increases.}
%   \label{fig:edm_ode_10_appendix}
% \end{figure}

% \begin{figure}[H]
%   \centering
%   \includegraphics[width=0.82\linewidth]{fig_transport_scenarios_grid_tarchic_vector_field.pdf}
%   \caption{Tarchic vector-field rendering of the synthetic transport scenarios. This figure records the learned-field geometry used in the two-dimensional pre-flight checks.}
%   \label{fig:tarchic_vector_field_appendix}
% \end{figure}

\begin{figure}[H]
  \centering
  \includegraphics[width=0.82\linewidth]{fig_transport_scenarios_grid_tarchic_vector_field_rose_points.pdf}
  \caption{Synthetic transport scenarios used as controlled probes. The grid contrasts continuous--continuous, continuous--discrete, discrete--continuous, and discrete--discrete endpoint structure. The point of this panel is to show that the bridge path can face different endpoint constraints even before any high-dimensional neural sampler is introduced..}
  \label{fig:tarchic_rose_points_appendix}
\end{figure}

% \begin{figure}[H]
%   \centering
%   \includegraphics[width=0.82\linewidth]{fig_transport_scenarios_grid_tarchic_vector_field_rose_points_alt_dc.pdf}
%   \caption{Alternative discrete--continuous rose-point transport view. This panel supports the claim that schedule effects depend on endpoint geometry.}
%   \label{fig:tarchic_alt_dc_appendix}
% \end{figure}

% \begin{figure}[H]
%   \centering
%   \includegraphics[width=0.82\linewidth]{fig_transport_scenarios_grid_tarchic_vector_field_rose_points_alt_dc_no_speedbar.pdf}
%   \caption{Alternative discrete--continuous rose-point view without the speedbar. This version is useful when the vector-field geometry should be inspected without color-scale emphasis.}
%   \label{fig:tarchic_alt_dc_no_speedbar_appendix}
% \end{figure}

% \begin{figure}[H]
%   \centering
%   \includegraphics[width=0.82\linewidth]{fig_transport_scenarios_grid_tarchic_vector_field_rose_points_dc_gaussian_modes_no_speedbar.pdf}
%   \caption{Discrete--continuous Gaussian-mode transport view without the speedbar. This is an additional controlled geometry used to inspect learned bridge behavior.}
%   \label{fig:tarchic_dc_gaussian_modes_appendix}
% \end{figure}

% \begin{figure}[H]
%   \centering
%   \includegraphics[width=0.82\linewidth]{fig_transport_scenarios_grid_tarchic_vector_field_rose_points_dc_no_cheat_flow_no_speedbar.pdf}
%   \caption{Discrete--continuous no-cheat-flow transport view. This panel is included as a diagnostic for learned-field structure, not as high-dimensional evidence.}
%   \label{fig:tarchic_dc_no_cheat_appendix}
% \end{figure}

% \begin{figure}[H]
%   \centering
%   \includegraphics[width=0.82\linewidth]{fig_transport_scenarios_grid_tarchic_vector_field_rose_points_learned_ode_no_speedbar.pdf}
%   \caption{Learned ODE vector-field view for the rose-point transports. This is the closest visual companion to the two-dimensional learned-field pre-flight experiment.}
%   \label{fig:tarchic_learned_ode_appendix}
% \end{figure}

% \begin{figure}[H]
%   \centering
%   \begin{subfigure}[t]{0.49\linewidth}
%     \centering
%     \includegraphics[width=\linewidth]{figures__EDM_trajectory_grid_ode_heun_25steps_linear.pdf}
%     \caption{Linear.}
%   \end{subfigure}
%   \hfill
%   \begin{subfigure}[t]{0.49\linewidth}
%     \centering
%     \includegraphics[width=\linewidth]{figures__EDM_trajectory_grid_ode_heun_25steps_entropic_log1p.pdf}
%     \caption{Cond--marg log1p.}
%   \end{subfigure}
%   \caption{EDM ODE-Heun trajectory grids at 25 steps. The panel shows how schedule differences become less pronounced as NFE increases.}
%   \label{fig:edm_ode_25_linear_log1p_appendix}
% \end{figure}

% \begin{figure}[H]
%   \centering
%   \includegraphics[width=0.72\linewidth]{alphaflow_condmarg_plddt_mean.pdf}
%   \caption{AlphaFlow mean pLDDT under cond--marg scheduling. This is endpoint-confidence evidence and should be read with \Cref{app:plddt_metric}.}
%   \label{fig:alphaflow_condmarg_plddt_mean_appendix}
% \end{figure}

\begin{figure}[H]
  \centering
  \includegraphics[width=0.78\linewidth]{cameo_alphaflow_condmarg_trajectory_plddt_mean_all_frames_by_dataset_size.pdf}
  \caption{CAMEO AlphaFlow trajectory pLDDT by dataset size. This panel documents endpoint-confidence trends across the trajectory and does not replace the cond--marg schedule diagnostic.}
  \label{fig:cameo_alphaflow_plddt_appendix}
\end{figure}

\begin{figure}[H]
  \centering
  \includegraphics[width=0.98\linewidth]{alphaflow_protein_testing_sankey.pdf}
  \caption{AlphaFlow protein-testing flow summary. This figure records the filtering and evaluation path behind the protein evidence.}
  \label{fig:alphaflow_testing_sankey_appendix}
\end{figure}

% \begin{figure}[H]
%   \centering
%   \includegraphics[width=0.72\linewidth]{\detokenize{plddt Distribution.pdf}}
%   \caption{Distribution of pLDDT endpoint-confidence scores. The distribution motivates treating pLDDT as a useful but limited proxy rather than as direct evidence of schedule geometry.}
%   \label{fig:plddt_distribution_appendix}
% \end{figure}
% To complement the FID table and trajectory grids, we include a dynamic-time-warping (DTW) diagnostic between scheduler time series. This is schedule-geometry evidence, not endpoint-quality evidence: it measures how much index warping is needed to align one scheduler with a cond--marg reference. At 10 steps, cond--marg \(\log(1+r)\) remains close to linear but distinguishable, whereas raw cond--marg allocation requires stronger warping. This supports using log tempering as a stable version of the endpoint-heavy entropy signal.

% \begin{figure}[H]
%   \centering
%   \includegraphics[width=\linewidth]{alphaflow_schedule_dtw_warping_paths_10step.pdf}
%   \caption{AlphaFlow ten-step dynamic-time-warping alignment paths between scheduler time series. Columns are query schedulers; rows use raw cond--marg and log-tempered cond--marg schedules as AlphaFlow references. The dark path is the optimal DTW alignment, and the color scale gives local index-pair distance.}
%   \label{fig:schedule_dtw_warping_paths_appendix}
% \end{figure}

\begin{figure}[H]
  \centering
  \includegraphics[width=0.82\linewidth]{figures__EDM_trajectory_grid_ode_heun_25steps_entropic_log1p.pdf}
  \caption{EDM ODE-Heun trajectory grid at 25 steps. The reader should observe that well-placed schedules begin to converge in visual behavior, matching the quantitative convergence in FID.}
  \label{fig:edm_ode_25_appendix}
\end{figure}

% \begin{figure}[H]
%   \centering
%   \includegraphics[width=0.82\linewidth]{EDM_trajectory_grid_sde_heun_5steps.pdf}
%   \caption{EDM SDE-Heun trajectory grid at 5 steps. This panel illustrates why the same time allocation can matter differently under stochastic sampling: injected noise can partially smooth grid-placement errors that are sharper in deterministic ODE integration.}
%   \label{fig:edm_sde_5_appendix}
% \end{figure}

% \begin{figure}[H]
%   \centering
%   \includegraphics[width=0.82\linewidth]{AlphaFlow_trajectory_grid_5steps.pdf}
%   \caption{AlphaFlow trajectory grid at 5 steps. The plot is a coarse-budget visualization of the protein flow path and should be read as a grid diagnostic rather than as a standalone structural-quality metric.}
%   \label{fig:alphaflow_grid_5_appendix}
% \end{figure}

% \begin{figure}[H]
%   \centering
%   \includegraphics[width=0.82\linewidth]{AlphaFlow_trajectory_grid_10steps.pdf}
%   \caption{AlphaFlow trajectory grid at 10 steps. The additional nodes make the transition path easier to inspect and provide context for the main-text 10-step protein rendering.}
%   \label{fig:alphaflow_grid_10_appendix}
% \end{figure}

% \begin{figure}[H]
%   \centering
%   \includegraphics[width=0.82\linewidth]{AlphaFlow_trajectory_grid_25steps.pdf}
%   \caption{AlphaFlow trajectory grid at 25 steps. This panel shows the higher-budget version of the same path visualization, where schedule-induced spacing is still visible but endpoint metrics become less sensitive to individual node placement.}
%   \label{fig:alphaflow_grid_25_appendix}
% \end{figure}

% \begin{figure}[H]
%   \centering
%   \begin{subfigure}[t]{0.48\linewidth}
%     \centering
%     \includegraphics[width=\linewidth]{AlphaFlow_entropic_cond_marg_protein_evolution_5steps_2d_7skj_A.pdf}
%     \caption{Five-step 2D evolution.}
%   \end{subfigure}
%   \hfill
%   \begin{subfigure}[t]{0.48\linewidth}
%     \centering
%     \includegraphics[width=\linewidth]{AlphaFlow_entropic_cond_marg_protein_evolution_10steps_3d_7skj_A.pdf}
%     \caption{Ten-step 3D evolution.}
%   \end{subfigure}
%   \caption{AlphaFlow conditional--marginal protein evolution for the same target under different visualizations. These panels support trajectory inspection and do not replace quantitative endpoint metrics.}
%   \label{fig:alphaflow_evolution_6p5x_appendix}
% \end{figure}

% \begin{figure}[H]
%   \centering
%   \includegraphics[width=0.72\linewidth]{AlphaFlow_entropic_cond_marg_protein_evolution_25steps_3d_7skj_A.pdf}
%   \caption{AlphaFlow 25-step 3D evolution for the same canonical protein target. This rendering is included to show that the schedule visualization is not tied to a single structure.}
%   \label{fig:alphaflow_evolution_7lp1_appendix}
% \end{figure}

\begin{figure}[H]
  \centering
  \begin{subfigure}[t]{0.48\linewidth}
    \centering
    \includegraphics[width=\linewidth]{AlphaFlow_entropic_cond_marg_protein_evolution_5steps_2d_7skj_A.pdf}
    \caption{2D view.}
  \end{subfigure}
  \hfill
  \begin{subfigure}[t]{0.48\linewidth}
    \centering
    \includegraphics[width=\linewidth]{AlphaFlow_entropic_cond_marg_protein_evolution_5steps_3d_7skj_A.pdf}
    \caption{3D view.}
  \end{subfigure}
  \caption{AlphaFlow cond--marg protein evolution at 5 steps for the same target. The paired 2D and 3D renderings are trajectory diagnostics, not endpoint-quality metrics.}
  \label{fig:alphaflow_evolution_5steps_appendix}
\end{figure}

\begin{figure}[H]
  \centering
  \begin{subfigure}[t]{0.48\linewidth}
    \centering
    \includegraphics[width=\linewidth]{AlphaFlow_entropic_cond_marg_protein_evolution_10steps_2d_7skj_A.pdf}
    \caption{2D view.}
  \end{subfigure}
  \hfill
  \begin{subfigure}[t]{0.48\linewidth}
    \centering
    \includegraphics[width=\linewidth]{AlphaFlow_entropic_cond_marg_protein_evolution_10steps_3d_7skj_A.pdf}
    \caption{3D view.}
  \end{subfigure}
  \caption{AlphaFlow cond--marg protein evolution at 10 steps for the same target. The paired views make the trajectory easier to inspect without changing the endpoint metric interpretation.}
  \label{fig:alphaflow_evolution_10steps_appendix}
\end{figure}

    % \begin{figure}[H]
    %   \centering
    %   \includegraphics[width=0.78\linewidth]{cond_marg_timeseries_10step_all_schedulers_edm_alphaflow.pdf}
    %   \caption{Ten-step scheduler time series for EDM and AlphaFlow. This appendix diagnostic is included for reference only; the main text keeps the simpler log1p-versus-linear comparison.}
    %   \label{fig:scheduler_shapes_appendix}
    % \end{figure}

\begin{figure}[H]
  \centering
  \begin{subfigure}[t]{0.48\linewidth}
    \centering
    \includegraphics[width=\linewidth]{AlphaFlow_entropic_cond_marg_protein_evolution_25steps_2d_7skj_A.pdf}
    \caption{2D view.}
  \end{subfigure}
  \hfill
  \begin{subfigure}[t]{0.48\linewidth}
    \centering
    \includegraphics[width=\linewidth]{AlphaFlow_entropic_cond_marg_protein_evolution_25steps_3d_7skj_A.pdf}
    \caption{3D view.}
  \end{subfigure}
  \caption{AlphaFlow cond--marg protein evolution at 25 steps for the same target. This higher-NFE view supports trajectory inspection and remains separate from quantitative pLDDT claims.}
  \label{fig:alphaflow_evolution_25steps_appendix}
\end{figure}

\section{Computational Notes}
\label{app:computational}
\paragraph{The inference data for our test runs was collected on a SLURM-managed HPC GPU cluster using NVIDIA H100 80GB HBM3 GPUs}. The main inference sweeps ran as exclusive 8-GPU node jobs, typically requesting gpus-per-node=8 with torchrun-nproc\_per\_node=8, about 80 requested CPU cores per job, 128 CPU cores allocated by SLURM, and roughly 2 TB of node memory. Our runs also included several smaller 1-GPU AlphaFlow precompute jobs with 16 CPUs and about 250 GB allocated memory. Counting allocated GPU wall time, including failed and timed-out retry attempts because they still consumed resources, and excluding CPU-only postprocessing, our runs consumed about 333.67 H100 GPU-hours in total: 183.52 GPU-hours for EDM and 150.15 GPU-hours for AlphaFlow, including experimentation. If the earlier invalid or cancelled launches are included as an all-in accounting number, the estimate rises slightly to about 334.42 H100 GPU-ours.

\paragraph{Computational cost.} Our method is computationally light relative to both EDM and AlphaFlow because it is a scheduler rather than a new generative backbone. At a fixed number of function evaluations (NFE), it uses the same trained-model calls as the corresponding baseline linear, cosine, sigmoid, or power schedule, and only changes the time grid on which those evaluations are placed. The additional online overhead consists of reading or interpolating a precomputed entropy-rate curve and constructing the schedule, which is $\mathcal{O}(K)$ for $K$ sampling steps and negligible compared with a single neural-network evaluation. The offline entropy-curve estimation incurs a one-time cost proportional to the number of time points and probe samples, but this cost is amortized across all subsequent samples and does not change the asymptotic inference complexity.

For EDM, inference is dominated by denoising-network evaluations. If $B$ is the batch size, $K$ is the number of model evaluations, and $C_{\mathrm{EDM}}$ is the cost of one EDM U-Net evaluation at the target image resolution, then sampling costs approximately
\[
\mathcal{O}(B K C_{\mathrm{EDM}}),
\]
up to constant-factor differences between Euler, Heun, and SDE-corrector implementations. Our entropic schedule has the same complexity at matched NFE; any improvement comes from allocating the same evaluations more effectively, or from reaching a target quality with fewer evaluations.

For AlphaFlow, the dominant cost is substantially larger because each denoising or integration step invokes a protein-structure model, here the ESMFold-based AlphaFlow checkpoint. For a protein of length $L$, each model call scales at least quadratically in $L$ due to attention and pairwise geometric representations, with additional architecture-dependent costs in the structure module. AlphaFlow sampling therefore scales roughly as
\[
\mathcal{O}(B K C_{\mathrm{AF}}(L)),
\]
where $C_{\mathrm{AF}}(L)$ grows steeply with protein length and memory also increases strongly with $L$. Our scheduler does not introduce additional AlphaFlow passes; it only changes the temporal allocation of the same 5-, 10-, or 25-step budgets.

%%%%%%%%%%%%%%%%%%%%%%%%%%%%%%%%%%%%%%%%%%%%%%%%%%%%%%%%%%%%

% \newpage
% \input{checklist.tex}


\end{document}
